\probchapter{第二章\quad 一元函数微分学}

\probsec{§1. 显函数的导数}

\begin{prob}[S]{821}
若 $x$ 由 $1$ 变到 $1000$, 求自变量 $x$ 的增量 $\Delta x$ 和函数 $y = \lg x$ 的相应增量 $\Delta y$.
\end{prob}

\begin{prob}[S]{822}
若 $x$ 由 $0.01$ 变到 $0.001$, 求自变量 $x$ 的增量 $\Delta x$ 和函数 $y = \frac{1}{x^2}$ 的相应增量 $\Delta y$.
\end{prob}

\begin{prob}[L]{823}
设:
\begin{enumerate}[label=(\alph*), itemsep=2pt, topsep=3pt]
\item $y = ax + b$;
\item $y = ax^2 + bx + c$;
\item $y = a^x$.
\end{enumerate}
若变量 $x$ 的增量为 $\Delta x$, 求增量 $\Delta y$.
\end{prob}

\begin{prob}[M]{824}
证明:
\begin{enumerate}[label=(\alph*), itemsep=2pt, topsep=3pt]
\item $\Delta[f(x) + g(x)] = \Delta f(x) + \Delta g(x)$;
\item $\Delta[f(x)g(x)] = g(x + \Delta x)\Delta f(x) + f(x)\Delta g(x)$.
\end{enumerate}
\end{prob}

\begin{prob}[L]{825}
过曲线 $y = x^2$ 上的两点 $A(2, 4)$ 和 $A'(2 + \Delta x, 4 + \Delta y)$ 引割线 $AA'$, 求此割线的斜率, 设:
\begin{enumerate}[label=(\alph*), itemsep=2pt, topsep=3pt]
\item $\Delta x = 1$;
\item $\Delta x = 0.1$;
\item $\Delta x = 0.01$;
\item $\Delta x$ 为任意小量.
\end{enumerate}
在该曲线上 $A$ 点的切线的斜率等于什么?
\end{prob}

\begin{prob}[L]{826}
利用函数 $y = x^3$ 把 $Ox$ 轴上的线段 $1 \leqslant x \leqslant 1 + h$ 映射到 $Oy$ 轴上, 求其平均伸长系数. 设:
\begin{enumerate}[label=(\alph*), itemsep=2pt, topsep=3pt]
\item $h = 0.1$;
\item $h = 0.01$;
\item $h = 0.001$,
\end{enumerate}
计算此系数的值. 在 $x = 1$, 此映射的伸长系数等于什么?
\end{prob}

\begin{prob}[L]{827}
动点沿 $Ox$ 轴运动的规律由下式给出:
\[
x = 10t + 5t^2,
\]
式中 $t$ 为以 $\mathrm{s}$ (秒) 计的时间, $x$ 为以 $\mathrm{m}$ (米) 计的距离. 求在 $20 \leqslant t \leqslant 20 + \Delta t$ 时间内运动的平均速度. 设:
\begin{enumerate}[label=(\alph*), itemsep=2pt, topsep=3pt]
\item $\Delta t = 1$;
\item $\Delta t = 0.1$;
\item $\Delta t = 0.01$,
\end{enumerate}
计算此速度的值. 当 $t = 20$ 时运动速度等于什么?
\end{prob}

\begin{prob}[L]{828}
根据导数的定义, 直接求下列函数的导数:
\begin{enumerate}[label=(\alph*), itemsep=2pt, topsep=3pt]
\item $x^2$;
\item $x^3$;
\item $\frac{1}{x}$;
\item $\sqrt{x}$;
\item $\sqrt[3]{x}$;
\item $\tan x$;
\item $\cot x$;
\item $\arcsin x$;
\item $\arccos x$;
\item $\arctan x$.
\end{enumerate}
\end{prob}

\begin{prob}[M]{829}
设 $f(x) = (x - 1)(x - 2)^2(x - 3)^3$, 求 $f'(1)$, $f'(2)$ 和 $f'(3)$.
\end{prob}

\begin{prob}[S]{830}
设 $f(x) = x^2\sin(x - 2)$, 求 $f'(2)$.
\end{prob}

\begin{prob}[M]{831}
设 $f(x) = x + (x - 1)\arcsin\sqrt{\frac{x}{x + 1}}$, 求 $f'(1)$.
\end{prob}

\begin{prob}[M]{832}
设函数 $f(x)$ 在 $a$ 点可微, 求
\[
\lim_{x \to a} \frac{f(x) - f(a)}{x - a}.
\]
\end{prob}

\begin{prob}[L]{833}
证明: 若函数 $f(x)$ 可微, $n$ 为正整数, 则
\[
\lim_{n \to \infty} n\left[ f\left(x + \frac{1}{n}\right) - f(x) \right] = f'(x). \tag{1}
\]
反之, 若对于函数 $f(x)$ 有极限 $(1)$ 存在, 则可否断定此函数有导数? 研究狄利克雷函数的例子 (参阅习题 734).
\end{prob}

\probgroup{利用导数表, 求下列函数的导数:}

\begin{prob}[S]{834}
$y = 2 + x - x^2$. 问 $y'(0)$, $y'\left(\frac{1}{2}\right)$, $y'(1)$, $y'(-10)$ 等于什么?
\end{prob}

\begin{prob}[L]{835}
$y = \frac{x^3}{3} + \frac{x^2}{2} - 2x$. 当 $x$ 为何值时:
\begin{enumerate}[label=(\alph*), itemsep=2pt, topsep=3pt]
\item $y'(x) = 0$;
\item $y'(x) = -2$;
\item $y'(x) = 10$?
\end{enumerate}
\end{prob}

\begin{prob}[S]{836}
$y = a^5 + 5a^3x^2 - x^5$.
\end{prob}

\begin{prob}[S]{837}
$y = \frac{ax + b}{a + b}$.
\end{prob}

\begin{prob}[S]{838}
$y = (x - a)(x - b)$.
\end{prob}

\begin{prob}[S]{839}
$y = (x + 1)(x + 2)^2(x + 3)^3$.
\end{prob}

\begin{prob}[S]{840}
$y = (x\sin\alpha + \cos\alpha)(x\cos\alpha - \sin\alpha)$.
\end{prob}

\begin{prob}[S]{841}
$y = (1 + nx^m)(1 + mx^n)$.
\end{prob}

\begin{prob}[M]{842}
(a) $y = (1 - x)(1 - x^2)^2(1 - x^3)^3$; \qquad (b) $y = (5 + 2x)^{10}(3 - 4x)^{20}$.
\end{prob}

\begin{prob}[S]{843}
$y = \frac{1}{x} + \frac{2}{x^2} + \frac{3}{x^3}$.
\end{prob}

\begin{prob}[M]{844}
证明: 公式 $\left(\frac{ax + b}{cx + d}\right)' = \frac{\begin{vmatrix} a & b \\ c & d \end{vmatrix}}{(cx + d)^2}$.
\end{prob}

\probgroup{求下列函数的导数:}

\begin{prob}[S]{845}
$y = \frac{2x}{1 - x^2}$.
\end{prob}

\begin{prob}[S]{846}
$y = \frac{1 + x - x^2}{1 - x + x^2}$.
\end{prob}

\begin{prob}[S]{847}
$y = \frac{x}{(1 - x)^2(1 + x)^3}$.
\end{prob}

\begin{prob}[S]{848}
$y = \frac{(2 - x^2)(3 - x^3)}{(1 - x)^2}$.
\end{prob}

\begin{prob}[S]{849}
$y = \frac{(1 - x)^p}{(1 + x)^q}$.
\end{prob}

\begin{prob}[S]{850}
$y = \frac{x^p(1 - x)^q}{1 + x}$.
\end{prob}

\begin{prob}[S]{851}
$y = x + \sqrt{x} + \sqrt[3]{x}$.
\end{prob}

\begin{prob}[S]{852}
$y = \frac{1}{x} + \frac{1}{\sqrt{x}} + \frac{1}{\sqrt[3]{x}}$.
\end{prob}

\begin{prob}[S]{853}
$y = \sqrt[3]{x^2} - \frac{2}{\sqrt{x}}$.
\end{prob}

\begin{prob}[S]{854}
$y = x\sqrt{1 + x^2}$.
\end{prob}

\begin{prob}[S]{855}
$y = (1 + x)\sqrt{2 + x^2}\sqrt[3]{3 + x^3}$.
\end{prob}

\begin{prob}[S]{856}
$y = \sqrt[m+n]{(1 - x)^m(1 + x)^n}$.
\end{prob}

\begin{prob}[S]{857}
$y = \frac{x}{\sqrt{a^2 - x^2}}$.
\end{prob}

\begin{prob}[S]{858}
$y = \sqrt[3]{\frac{1 + x^3}{1 - x^3}}$.
\end{prob}

\begin{prob}[S]{859}
$y = \frac{1}{\sqrt{1 + x^2}(x + \sqrt{1 + x^2})}$.
\end{prob}

\begin{prob}[S]{860}
$y = \sqrt{x + \sqrt{x + \sqrt{x}}}$.
\end{prob}

\begin{prob}[S]{861}
$y = \sqrt[3]{1 + \sqrt[3]{1 + \sqrt[3]{x}}}$.
\end{prob}

\begin{prob}[S]{862}
$y = \cos 2x - 2\sin x$.
\end{prob}

\begin{prob}[S]{863}
$y = (2 - x^2)\cos x + 2x\sin x$.
\end{prob}

\begin{prob}[S]{864}
$y = \sin(\cos^2 x)\cdot\cos(\sin^2 x)$.
\end{prob}

\begin{prob}[S]{865}
$y = \sin^n x \cos nx$.
\end{prob}

\begin{prob}[S]{866}
$y = \sin[\sin(\sin x)]$.
\end{prob}

\begin{prob}[S]{867}
$y = \frac{\sin^2 x}{\sin x^2}$.
\end{prob}

\begin{prob}[S]{868}
$y = \frac{\cos x}{2\sin^2 x}$.
\end{prob}

\begin{prob}[S]{869}
$y = \frac{1}{\cos^n x}$.
\end{prob}

\begin{prob}[S]{870}
$y = \frac{\sin x - x\cos x}{\cos x + x\sin x}$.
\end{prob}

\begin{prob}[S]{871}
$y = \tan\frac{x}{2} - \cot\frac{x}{2}$.
\end{prob}

\begin{prob}[S]{872}
$y = \tan x - \frac{1}{3}\tan^3 x + \frac{1}{5}\tan^5 x$.
\end{prob}

\begin{prob}[S]{873}
$y = 4\sqrt[3]{\cot^2 x} + \sqrt[3]{\cot^8 x}$.
\end{prob}

\begin{prob}[S]{874}
$y = \sec^2\frac{x}{a} + \csc^2\frac{x}{a}$.
\end{prob}

\begin{prob}[S]{875}
$y = \sin[\cos^2(\tan^3 x)]$.
\end{prob}

\begin{prob}[S]{876}
$y = \mathrm{e}^{-x^2}$.
\end{prob}

\begin{prob}[S]{877}
$y = 2^{\tan\frac{1}{x}}$.
\end{prob}

\begin{prob}[S]{878}
$y = \mathrm{e}^x(x^2 - 2x + 2)$.
\end{prob}

\begin{prob}[S]{879}
$y = \left[\frac{1 - x^2}{2}\sin x - \frac{(1 + x)^2}{2}\cos x\right]\mathrm{e}^{-x}$.
\end{prob}

\begin{prob}[S]{880}
$y = \mathrm{e}^x\left(1 + \cot\frac{x}{2}\right)$.
\end{prob}

\begin{prob}[S]{881}
$y = \frac{\ln 3 \cdot \sin x + \cos x}{3^x}$.
\end{prob}

\begin{prob}[S]{882}
$y = \mathrm{e}^{ax}\frac{a\sin bx - b\cos bx}{\sqrt{a^2 + b^2}}$.
\end{prob}

\begin{prob}[S]{883}
$y = \mathrm{e}^x + \mathrm{e}^{\mathrm{e}^x} + \mathrm{e}^{\mathrm{e}^{\mathrm{e}^x}}$.
\end{prob}

\begin{prob}[S]{884}
$y = \left(\frac{a}{b}\right)^x \left(\frac{b}{x}\right)^a \left(\frac{x}{a}\right)^b$ $(a > 0, b > 0)$.
\end{prob}

\begin{prob}[S]{885}
$y = x^{a^a} + a^{x^a} + a^{a^x}$ $(a > 0)$.
\end{prob}

\begin{prob}[S]{886}
$y = \lg^3 x^2$.
\end{prob}

\begin{prob}[S]{887}
$y = \ln(\ln(\ln x))$.
\end{prob}

\begin{prob}[S]{888}
$y = \ln(\ln^2(\ln^3 x))$.
\end{prob}

\begin{prob}[S]{889}
$y = \frac{1}{2}\ln(1 + x) - \frac{1}{4}\ln(1 + x^2) - \frac{1}{2(1 + x)}$.
\end{prob}

\begin{prob}[S]{890}
$y = \frac{1}{4}\ln\frac{x^2 - 1}{x^2 + 1}$.
\end{prob}

\begin{prob}[S]{891}
$y = \frac{1}{4(1 + x^4)} + \frac{1}{4}\ln\frac{x^4}{1 + x^4}$.
\end{prob}

\begin{prob}[S]{892}
$y = \frac{1}{2\sqrt{6}}\ln\frac{x\sqrt{3} - \sqrt{2}}{x\sqrt{3} + \sqrt{2}}$.
\end{prob}

\begin{prob}[S]{893}
$y = \frac{1}{1 - k}\ln\frac{1 + x}{1 - x} + \frac{\sqrt{k}}{1 - k}\ln\frac{1 + x\sqrt{k}}{1 - x\sqrt{k}}$ $(0 < k < 1)$.
\end{prob}

\begin{prob}[S]{894}
$y = \sqrt{x + 1} - \ln(1 + \sqrt{x + 1})$.
\end{prob}

\begin{prob}[S]{895}
$y = \ln(x + \sqrt{x^2 + 1})$.
\end{prob}

\begin{prob}[S]{896}
$y = x\ln(x + \sqrt{1 + x^2}) - \sqrt{1 + x^2}$.
\end{prob}

\begin{prob}[S]{897}
$y = x\ln^2(x + \sqrt{1 + x^2}) - 2\sqrt{1 + x^2}\ln(x + \sqrt{1 + x^2}) + 2x$.
\end{prob}

\begin{prob}[S]{898}
$y = \frac{x}{2}\sqrt{x^2 + a^2} + \frac{a^2}{2}\ln(x + \sqrt{x^2 + a^2})$.
\end{prob}

\begin{prob}[S]{899}
$y = \frac{1}{2\sqrt{ab}}\ln\frac{\sqrt{a} + x\sqrt{b}}{\sqrt{a} - x\sqrt{b}}$ $(a > 0, b > 0)$.
\end{prob}

\begin{prob}[S]{900}
$y = \frac{2 + 3x^2}{x^4}\sqrt{1 - x^2} + 3\ln\frac{1 + \sqrt{1 - x^2}}{x}$.
\end{prob}

\begin{prob}[S]{901}
$y = \ln\tan\frac{x}{2}$.
\end{prob}

\begin{prob}[S]{902}
$y = \ln\tan\left(\frac{x}{2} + \frac{\pi}{4}\right)$.
\end{prob}

\begin{prob}[S]{903}
$y = \frac{1}{2}\cot^2 x + \ln\sin x$.
\end{prob}

\begin{prob}[S]{904}
$y = \ln\sqrt{\frac{1 - \sin x}{1 + \sin x}}$.
\end{prob}

\begin{prob}[S]{905}
$y = -\frac{\cos x}{2\sin^2 x} + \ln\sqrt{\frac{1 + \cos x}{\sin x}}$.
\end{prob}

\begin{prob}[S]{906}
$y = \ln\frac{b + a\cos x + \sqrt{b^2 - a^2}\sin x}{a + b\cos x}$ $(0 \leqslant |a| < |b|)$.
\end{prob}

\begin{prob}[S]{907}
$y = \frac{1}{x}(\ln^3 x + 3\ln^2 x + 6\ln x + 6)$.
\end{prob}

\begin{prob}[S]{908}
$y = \frac{1}{4x^4}\ln\frac{1}{x} - \frac{1}{16x^4}$.
\end{prob}

\begin{prob}[S]{909}
$y = \frac{3}{2}\left(1 - \sqrt[3]{1 + x^2}\right)^2 + 3\ln\left(1 + \sqrt[3]{1 + x^2}\right)$.
\end{prob}

\begin{prob}[S]{910}
$y = \ln\left[\frac{1}{x} + \ln\left(\frac{1}{x} + \ln\frac{1}{x}\right)\right]$.
\end{prob}

\begin{prob}[S]{911}
$y = x[\sin(\ln x) - \cos(\ln x)]$.
\end{prob}

\begin{prob}[S]{912}
$y = \ln\tan\frac{x}{2} - \cos x \cdot \ln\tan x$.
\end{prob}

\begin{prob}[S]{913}
$y = \arcsin\frac{x}{2}$.
\end{prob}

\begin{prob}[S]{914}
$y = \arccos\frac{1 - x}{\sqrt{2}}$.
\end{prob}

\begin{prob}[S]{915}
$y = \arctan\frac{x^2}{a}$.
\end{prob}

\begin{prob}[S]{916}
$y = \frac{1}{\sqrt{2}}\operatorname{arccot}\frac{\sqrt{2}}{x}$.
\end{prob}

\begin{prob}[S]{917}
$y = \sqrt{x} - \arctan\sqrt{x}$.
\end{prob}

\begin{prob}[S]{918}
$y = x + \sqrt{1 - x^2}\cdot\arccos x$.
\end{prob}

\begin{prob}[S]{919}
$y = x\arcsin\sqrt{\frac{x}{1 + x}} + \arctan\sqrt{x} - \sqrt{x}$.
\end{prob}

\begin{prob}[S]{920}
$y = \arccos\frac{1}{x}$.
\end{prob}

\begin{prob}[S]{921}
$y = \arcsin(\sin x)$.
\end{prob}

\begin{prob}[S]{922}
$y = \arccos(\cos^2 x)$.
\end{prob}

\begin{prob}[S]{923}
$y = \arcsin(\sin x - \cos x)$.
\end{prob}

\begin{prob}[S]{924}
$y = \arccos\sqrt{1 - x^2}$.
\end{prob}

\begin{prob}[S]{925}
$y = \arctan\frac{1 + x}{1 - x}$.
\end{prob}

\begin{prob}[S]{926}
$y = \operatorname{arccot}\left(\frac{\sin x + \cos x}{\sin x - \cos x}\right)$.
\end{prob}

\begin{prob}[S]{927}
$y = \frac{2}{\sqrt{a^2 - b^2}} \arctan \left( \sqrt{\frac{a - b}{a + b}} \tan \frac{x}{2} \right) \ (a > b \geqslant 0)$.
\end{prob}

\begin{prob}[S]{928}
$y = \arcsin \frac{1 - x^2}{1 + x^2}$.
\end{prob}

\begin{prob}[S]{929}
$y = \frac{1}{\arccos^2(x^2)}$.
\end{prob}

\begin{prob}[S]{930}
$y = \arctan x + \frac{1}{3} \arctan(x^3)$.
\end{prob}

\begin{prob}[S]{931}
$y = \ln(1 + \sin^2 x) - 2 \sin x \cdot \arctan(\sin x)$.
\end{prob}

\begin{prob}[S]{932}
$y = \ln \left( \arccos \frac{1}{\sqrt{x}} \right)$.
\end{prob}

\begin{prob}[S]{933}
$y = \ln \frac{x + a}{\sqrt{x^2 + b^2}} + \frac{a}{b} \arctan \frac{x}{b}$.
\end{prob}

\begin{prob}[S]{934}
$y = \frac{x}{2} \sqrt{a^2 - x^2} + \frac{a^2}{2} \arcsin \frac{x}{a} \ (a > 0)$.
\end{prob}

\begin{prob}[S]{935}
$y = \frac{1}{6} \ln \frac{(x + 1)^2}{x^2 - x + 1} + \frac{1}{\sqrt{3}} \arctan \frac{2x - 1}{\sqrt{3}}$.
\end{prob}

\begin{prob}[S]{936}
$y = \frac{1}{4\sqrt{2}} \ln \frac{x^2 + x\sqrt{2} + 1}{x^2 - x\sqrt{2} + 1} - \frac{1}{2\sqrt{2}} \arctan \frac{x\sqrt{2}}{x^2 - 1}$.
\end{prob}

\begin{prob}[S]{937}
$y = x(\arcsin x)^2 + 2\sqrt{1 - x^2} \arcsin x - 2x$.
\end{prob}

\begin{prob}[S]{938}
$y = \frac{\arccos x}{x} + \frac{1}{2} \ln \frac{1 - \sqrt{1 - x^2}}{1 + \sqrt{1 - x^2}}$.
\end{prob}

\begin{prob}[S]{939}
$y = \arctan \sqrt{x^2 - 1} - \frac{\ln x}{\sqrt{x^2 - 1}}$.
\end{prob}

\begin{prob}[S]{940}
$y = \frac{\arcsin x}{\sqrt{1 - x^2}} + \frac{1}{2} \ln \frac{1 - x}{1 + x}$.
\end{prob}

\begin{prob}[S]{941}
$y = \frac{1}{12} \ln \frac{x^4 - x^2 + 1}{(x^2 + 1)^2} - \frac{1}{2\sqrt{3}} \arctan \frac{\sqrt{3}}{2x^2 - 1}$.
\end{prob}

\begin{prob}[S]{942}
$y = \frac{x^6}{1 + x^{12}} - \operatorname{arccot} x^6$.
\end{prob}

\begin{prob}[S]{943}
$y = \ln \frac{1 - \sqrt[3]{x}}{\sqrt{1 + \sqrt[3]{x} + \sqrt[3]{x^2}}} + \sqrt{3} \arctan \frac{1 + 2\sqrt[3]{x}}{\sqrt{3}}$.
\end{prob}

\begin{prob}[S]{944}
$y = \arctan \frac{x}{1 + \sqrt{1 - x^2}}$.
\end{prob}

\begin{prob}[S]{945}
$y = \operatorname{arccot} \frac{a - 2x}{2\sqrt{ax - x^2}} \ (a > 0)$.
\end{prob}

\begin{prob}[S]{946}
$y = \frac{3 - x}{2} \sqrt{1 - 2x - x^2} + 2\arcsin \frac{1 + x}{\sqrt{2}}$.
\end{prob}

\begin{prob}[S]{947}
$y = \frac{1}{4} \ln \frac{\sqrt[4]{1 + x^4} + x}{\sqrt[4]{1 + x^4} - x} - \frac{1}{2} \arctan \frac{\sqrt[4]{1 + x^4}}{x}$.
\end{prob}

\begin{prob}[S]{948}
$y = \arctan(\tan^2 x)$.
\end{prob}

\begin{prob}[S]{949}
$y = \sqrt{1 - x^2} \cdot \ln \sqrt{\frac{1 - x}{1 + x}} + \frac{1}{2} \ln \frac{1 - \sqrt{1 - x^2}}{1 + \sqrt{1 - x^2}} + \sqrt{1 - x^2} + \arcsin x$.
\end{prob}

\begin{prob}[S]{950}
$y = x\arctan x - \frac{1}{2} \ln(1 + x)^2 - \frac{1}{2} (\arctan x)^2$.
\end{prob}

\begin{prob}[S]{951}
$y = \ln(\mathrm{e}^x + \sqrt{1 + \mathrm{e}^{2x}})$.
\end{prob}

\begin{prob}[S]{952}
$y = \arctan(x + \sqrt{1 + x^2})$.
\end{prob}

\begin{prob}[S]{953}
$y = \arcsin \left( \frac{\sin a \sin x}{1 - \cos a \cos x} \right)$.
\end{prob}

\begin{prob}[S]{954}
$y = \frac{1}{4\sqrt{3}} \ln \frac{\sqrt{x^2 + 2} - x\sqrt{3}}{\sqrt{x^2 + 2} + x\sqrt{3}} + \frac{1}{2} \arctan \frac{\sqrt{x^2 + 2}}{x}$.
\end{prob}

\begin{prob}[S]{955}
$y = \frac{1}{2\sqrt{2}} \arctan \frac{x\sqrt{2}}{\sqrt{1 + x^4}} - \frac{1}{4\sqrt{2}} \ln \frac{\sqrt{1 + x^4} - x\sqrt{2}}{\sqrt{1 + x^4} + x\sqrt{2}}$.
\end{prob}

\begin{prob}[S]{956}
$y = \frac{x\sqrt{1 - x^2}}{1 + x^2} - \frac{3}{\sqrt{2}} \operatorname{arccot} \frac{x\sqrt{2}}{\sqrt{1 - x^2}}$.
\end{prob}

\begin{prob}[S]{957}
$y = \arccos(\sin x^2 - \cos x^2)$.
\end{prob}

\begin{prob}[S]{958}
$y = \arcsin(\sin x^2) + \arccos(\cos x^2)$.
\end{prob}

\begin{prob}[S]{959}
$y = \mathrm{e}^{m\arcsin x}[\cos(m\arcsin x) + \sin(m\arcsin x)]$.
\end{prob}

\begin{prob}[L]{960}
\begin{enumerate}[label=(\alph*), itemsep=2pt, topsep=3pt]
\item $y = \arctan \mathrm{e}^x - \ln \sqrt{\frac{\mathrm{e}^{2x}}{\mathrm{e}^{2x} + 1}}$;
\item $y = \sqrt{1 + \sqrt[3]{1 + \sqrt[4]{1 + x^4}}}$;
\item $y = \operatorname{arccot} \frac{1}{\sqrt{\cot \frac{1}{x^2}}}$;
\item $y = \ln^2(\sec 2^{\sqrt[3]{x}})$.
\end{enumerate}
\end{prob}

\begin{prob}[S]{961}
$y = x + x^x + x^{x^x} \ (x > 0)$.
\end{prob}

\begin{prob}[S]{962}
$y = x^{x^a} + x^{a^x} + a^{x^x} \ (a > 0, x > 0)$.
\end{prob}

\begin{prob}[S]{963}
$y = \sqrt[x]{x} \ (x > 0)$.
\end{prob}

\begin{prob}[S]{964}
$y = (\sin x)^{\cos x} + (\cos x)^{\sin x}$.
\end{prob}

\begin{prob}[S]{965.1}
$y = \frac{(\ln x)^x}{x^{\ln x}}$.
\end{prob}

\begin{prob}[S]{965.2}
$y = \left[ \frac{\arcsin(\sin^2 x)}{\arccos(\cos^2 x)} \right]^{\arctan^2 x}$.
\end{prob}

\begin{prob}[S]{966}
$y = \log_x \mathrm{e}$.
\end{prob}

\begin{prob}[S]{967}
$y = \ln(\cosh x) + \frac{1}{2\cosh^2 x}$.
\end{prob}

\begin{prob}[S]{968}
$y = \frac{\cosh x}{\sinh^2 x} - \ln\left(\coth \frac{x}{2}\right)$.
\end{prob}

\begin{prob}[S]{969}
$y = \arctan(\tanh x)$.
\end{prob}

\begin{prob}[S]{970}
$y = \arccos\left(\frac{1}{\cosh x}\right)$.
\end{prob}

\begin{prob}[S]{971}
$y = \frac{b}{a} x + \frac{2\sqrt{a^2 - b^2}}{a} \arctan \left( \sqrt{\frac{a - b}{a + b}} \tanh \frac{x}{2} \right) \ (0 \leqslant |b| < a)$.
\end{prob}

\begin{prob}[M]{972}
引入中间变量 $u = \cos^2 x$，求函数
\[
y = \ln(\cos^2 x + \sqrt{1 + \cos^4 x})
\]
的导数。
\end{prob}

\probgroup{利用 972 题所示的方法，求下列函数的导数：}

\begin{prob}[S]{973}
$y = (\arccos x)^2 \left[ \ln^2(\arccos x) - \ln(\arccos x) + \frac{1}{2} \right]$.
\end{prob}

\begin{prob}[S]{974}
$y = \frac{1}{2} \arctan(\sqrt[4]{1 + x^4}) + \frac{1}{4} \ln \frac{\sqrt[4]{1 + x^4} + 1}{\sqrt[4]{1 + x^4} - 1}$.
\end{prob}

\begin{prob}[S]{975}
$y = \frac{\mathrm{e}^{-x^2} \arcsin(\mathrm{e}^{-x^2})}{\sqrt{1 - \mathrm{e}^{-2x^2}}} + \frac{1}{2} \ln(1 - \mathrm{e}^{-2x^2})$.
\end{prob}

\begin{prob}[S]{976}
$y = \frac{a^x}{1 + a^{2x}} - \frac{1 - a^{2x}}{1 + a^{2x}} \operatorname{arccot} a^{-x}$.
\end{prob}

\begin{prob}[L]{977}
求下列函数的导数并作出函数及其导数的图像：
\begin{enumerate}[label=(\alph*), itemsep=2pt, topsep=3pt]
\item $y = |x|$;
\item $y = x|x|$;
\item $y = \ln |x|$.
\end{enumerate}
\end{prob}

\begin{prob}[L]{978}
求下列函数的导数：
\begin{enumerate}[label=(\alph*), itemsep=2pt, topsep=3pt]
\item $y = |(x - 1)^2 (x + 1)^3|$;
\item $y = |\sin^3 x|$;
\item $y = \arccos \frac{1}{|x|}$;
\item $y = [x] \sin^2 \pi x$.
\end{enumerate}
\end{prob}

\probgroup{求下列函数的导数并作出函数及其导数的图像：}

\begin{prob}[M]{979}
\[
y = \begin{cases}
1 - x, & -\infty < x < 1, \\
(1 - x)(2 - x), & 1 \leqslant x \leqslant 2, \\
-(2 - x), & 2 < x < +\infty.
\end{cases}
\]
\end{prob}

\begin{prob}[M]{980}
\[
y = \begin{cases}
(x - a)^2 (x - b)^2, & a \leqslant x \leqslant b, \\
0, & \text{在线段 } [a, b] \text{ 之外}.
\end{cases}
\]
\end{prob}

\begin{prob}[M]{981}
\[
y = \begin{cases}
x, & x < 0, \\
\ln(1 + x), & x \geqslant 0.
\end{cases}
\]
\end{prob}

\begin{prob}[M]{982}
\[
y = \begin{cases}
\arctan x, & |x| \leqslant 1, \\
\dfrac{\pi}{4} \operatorname{sgn} x + \dfrac{x - 1}{2}, & |x| > 1.
\end{cases}
\]
\end{prob}

\begin{prob}[M]{983}
\[
y = \begin{cases}
x^2 \mathrm{e}^{-x^2}, & |x| \leqslant 1, \\
\dfrac{1}{\mathrm{e}}, & |x| > 1.
\end{cases}
\]
\end{prob}

\begin{prob}[L]{984}
所给函数 $y = f(x)$ 的对数的导数称为此函数的对数导数：
\[
\frac{y'}{y} = \frac{\mathrm{d}}{\mathrm{d}x} \ln f(x) = \frac{f'(x)}{f(x)} \quad (f(x) > 0).
\]
求下列函数 $y$ 的对数导数：
\begin{enumerate}[label=(\alph*), itemsep=2pt, topsep=3pt]
\item $y = x \sqrt{\frac{1 - x}{1 + x}}$;
\item $y = \frac{x^2}{1 - x} \sqrt[3]{\frac{3 - x}{(3 + x)^2}}$;
\item $y = (x - a_1)^{\alpha_1} (x - a_2)^{\alpha_2} \cdots (x - a_n)^{\alpha_n}$;
\item $y = \left(x + \sqrt{1 + x^2}\right)^n$.
\end{enumerate}
\end{prob}

\begin{prob}[L]{985}
设 $\varphi(x)$ 及 $\psi(x)$ 为 $x$ 的可微函数. 求下列函数 $y$ 的导数：
\begin{enumerate}[label=(\alph*), itemsep=2pt, topsep=3pt]
\item $y = \sqrt{\varphi^2(x) + \psi^2(x)}$;
\item $y = \arctan \frac{\varphi(x)}{\psi(x)}$;
\item $y = \sqrt[\varphi(x)]{\psi(x)} \ (\varphi(x) \neq 0, \ \psi(x) > 0)$;
\item $y = \log_{\varphi(x)} \psi(x) \ (\varphi(x) > 0, \ \psi(x) > 0)$.
\end{enumerate}
\end{prob}

\begin{prob}[L]{986.1}
求 $y'$，设
\begin{enumerate}[label=(\alph*), itemsep=2pt, topsep=3pt]
\item $y = f(x^2)$;
\item $y = f(\sin^2 x) + f(\cos^2 x)$;
\item $y = f(\mathrm{e}^x) \cdot \mathrm{e}^{f(x)}$;
\item $y = f\{f[f(x)]\}$，
\end{enumerate}
其中 $f(u)$ 为可微函数.
\end{prob}

\begin{prob}[M]{986.2}
设 $f(x) = x(x - 1)(x - 2)\cdots(x - 1000)$，求 $f'(0)$.
\end{prob}

\begin{prob}[L]{987}
证明：$n$ 阶行列式微分法：
\[
\begin{vmatrix}
f_{11}(x) & f_{12}(x) & \cdots & f_{1n}(x) \\
\vdots & \vdots & & \vdots \\
f_{k1}(x) & f_{k2}(x) & \cdots & f_{kn}(x) \\
\vdots & \vdots & & \vdots \\
f_{n1}(x) & f_{n2}(x) & \cdots & f_{nn}(x)
\end{vmatrix}'
= \sum_{k=1}^{n}
\begin{vmatrix}
f_{11}(x) & f_{12}(x) & \cdots & f_{1n}(x) \\
\vdots & \vdots & & \vdots \\
f'_{k1}(x) & f'_{k2}(x) & \cdots & f'_{kn}(x) \\
\vdots & \vdots & & \vdots \\
f_{n1}(x) & f_{n2}(x) & \cdots & f_{nn}(x)
\end{vmatrix}.
\]
\end{prob}

\begin{prob}[M]{988}
设
\[
F(x) = \begin{vmatrix}
x - 1 & 1 & 2 \\
-3 & x & 3 \\
-2 & -3 & x + 1
\end{vmatrix},
\]
求 $F'(x)$.
\end{prob}

\begin{prob}[M]{989}
设
\[
F(x) = \begin{vmatrix}
x & x^2 & x^3 \\
1 & 2x & 3x^2 \\
0 & 2 & 6x
\end{vmatrix},
\]
求 $F'(x)$.
\end{prob}

\begin{prob}[M]{990}
已知函数的图像，近似地作出其导数的图像.
\end{prob}

\begin{prob}[M]{991}
证明：函数
\[
f(x) = \begin{cases}
x^2 \sin \dfrac{1}{x}, & x \neq 0, \\
0, & x = 0
\end{cases}
\]
有不连续的导数.
\end{prob}

\begin{prob}[L]{992}
在什么条件下，函数
\[
f(x) = \begin{cases}
x^n \sin \dfrac{1}{x}, & x \neq 0, \\
0, & x = 0
\end{cases}
\]
\begin{enumerate}[label=(\alph*), itemsep=2pt, topsep=3pt]
\item 在 $x = 0$ 处是连续的；
\item 在 $x = 0$ 处可微；
\item 在 $x = 0$ 处有连续的导数？
\end{enumerate}
\end{prob}

\begin{prob}[M]{993}
在什么条件下，函数
\[
f(x) = \begin{cases}
|x|^n \sin \dfrac{1}{|x|^m}, & x \neq 0, \\
0, & x = 0
\end{cases}
\quad (m > 0)
\]
(a) 于坐标原点的邻域上有有界的导数； \qquad (b) 在此邻域上有无穷导数？
\end{prob}

\begin{prob}[M]{994}
设
\[
f(x) = (x - a)\varphi(x),
\]
其中函数 $\varphi(x)$ 在 $x = a$ 处是连续的，求 $f'(a)$.
\end{prob}

\begin{prob}[M]{995}
设
\[
f(x) = |x - a|\varphi(x),
\]
其中 $\varphi(x)$ 为连续函数且 $\varphi(a) \neq 0$，证明：此函数在 $a$ 点没有导数. 单侧导数 $f'_-(a)$ 及 $f'_+(a)$ 等于什么？
\end{prob}

\begin{prob}[M]{996}
举出在已知点 $a_1, a_2, \cdots, a_n$ 没有导数的连续函数的例子.
\end{prob}

\begin{prob}[M]{997}
证明：函数
\[
f(x) = \begin{cases}
x^2 \left|\cos \dfrac{\pi}{x}\right|, & x \neq 0, \\
0, & x = 0
\end{cases}
\]
在点 $x = 0$ 的任何邻域上都有不可微点，但在该点是可微的. 作出此函数的略图.
\end{prob}

\begin{prob}[M]{998}
证明：函数
\[
f(x) = \begin{cases}
x^2, & x \text{ 为有理数}, \\
0, & x \text{ 为无理数}
\end{cases}
\]
仅在 $x = 0$ 时有导数.
\end{prob}

\begin{prob}[L]{999}
研究下列函数的可微性：
\begin{enumerate}[label=(\alph*), itemsep=2pt, topsep=3pt]
\item $y = |(x - 1)(x - 2)^2(x - 3)^3|$;
\item $y = |\cos x|$;
\item $y = |\pi^2 - x^2|\sin^2 x$;
\item $y = \arcsin(\cos x)$;
\item $y = \begin{cases} \dfrac{x - 1}{4}(x + 1)^2, & |x| \leqslant 1, \\ |x| - 1, & |x| > 1. \end{cases}$
\end{enumerate}
\end{prob}

\probgroup{求函数 $f(x)$ 的左导数 $f'_-(x)$ 和右导数 $f'_+(x)$，设：}

\begin{prob}[S]{1000}
$f(x) = |x|$.
\end{prob}

\begin{prob}[S]{1001}
$f(x) = [x] \sin \pi x$.
\end{prob}

\begin{prob}[M]{1002}
$f(x) = \begin{cases}
x \left|\cos \dfrac{\pi}{x}\right|, & x \neq 0, \\
0, & x = 0.
\end{cases}$
\end{prob}

\begin{prob}[S]{1003}
$f(x) = \sqrt{\sin x^2}$.
\end{prob}

\begin{prob}[M]{1004}
$f(x) = \begin{cases}
\dfrac{x}{1 + \mathrm{e}^{\frac{1}{x}}}, & x \neq 0, \\
0, & x = 0.
\end{cases}$
\end{prob}

\begin{prob}[S]{1005}
$f(x) = \sqrt{1 - \mathrm{e}^{-x^2}}$.
\end{prob}

\begin{prob}[S]{1006}
$f(x) = |\ln|x|| \ (x \neq 0)$.
\end{prob}

\begin{prob}[S]{1007}
$f(x) = \arcsin \frac{2x}{1 + x^2}$.
\end{prob}

\begin{prob}[M]{1008}
$f(x) = \begin{cases}
(x - 2)\arctan \dfrac{1}{x - 2}, & x \neq 2, \\
0, & x = 2.
\end{cases}$
\end{prob}

\begin{prob}[M]{1009.1}
证明：函数
\[
f(x) = \begin{cases}
x \sin \dfrac{1}{x}, & x \neq 0, \\
0, & x = 0
\end{cases}
\]
在点 $x = 0$ 连续，但在此点既无左导数，又无右导数.
\end{prob}

\begin{prob}[M]{1009.2}
设 $x_0$ 为函数 $f(x)$ 的第一类间断点. 表达式
\[
f'_-(x_0) = \lim_{h \to -0} \frac{f(x_0 + h) - f(x_0 - 0)}{h},
\]
\[
f'_+(x_0) = \lim_{h \to +0} \frac{f(x_0 + h) - f(x_0 + 0)}{h}
\]
称为函数 $f(x)$ 在点 $x_0$ 的广义单侧导数（左导数和右导数）.
求函数 $f(x)$ 在间断点 $x_0$ 的 $f'_-(x_0)$ 和 $f'_+(x_0)$：
\begin{enumerate}[label=(\alph*), itemsep=2pt, topsep=3pt]
  \item $f(x) = \dfrac{\sqrt{x^2 + x^3}}{x}$；
  \item $f(x) = \arctan \dfrac{1+x}{1-x}$；
  \item $f(x) = \dfrac{1}{1 + \mathrm{e}^{\frac{1}{x}}}$．
\end{enumerate}
\end{prob}

\begin{prob}[M]{1010}
设
\[
f(x) = \begin{cases}
x^2, & x \leqslant x_0, \\
ax + b, & x > x_0.
\end{cases}
\]
为了使函数 $f(x)$ 在点 $x = x_0$ 处连续而且可微，应当如何选取系数 $a$ 和 $b$？
\end{prob}

\begin{prob}[M]{1011}
设
\[
F(x) = \begin{cases}
f(x), & x \leqslant x_0, \\
ax + b, & x > x_0,
\end{cases}
\]
其中函数 $f(x)$ 在点 $x = x_0$ 处为左可微的．应当如何选取系数 $a$ 和 $b$，使函数 $F(x)$ 在点 $x_0$ 处连续而且可微？
\end{prob}

\begin{prob}[M]{1012}
适当地选定参数 $A$ 和 $c$，用立方抛物线
\[
y = A(x - a)(x - b)(x - c)
\]
在闭区间 $a \leqslant x \leqslant b$ 上把两条射线：
\[
y = k_1(x - a) \quad (-\infty < x < a) \quad \text{及} \quad y = k_2(x - b) \quad (b < x < +\infty)
\]
光滑地连接起来．
\end{prob}

\begin{prob}[M]{1013}
用抛物线 $y = a + bx^2$ ($|x| \leqslant c$)（其中 $a$ 与 $b$ 为未知的参数）去补充曲线 $y = \dfrac{m^2}{|x|}$ ($|x| > c$)，以便得到一条光滑曲线．
\end{prob}

\begin{prob}[M]{1014}
若 (a) 函数 $f(x)$ 在点 $x_0$ 有导数，而函数 $g(x)$ 在此点没有导数；(b) 函数 $f(x)$ 和 $g(x)$ 二者在点 $x_0$ 都没有导数，可否断定它们的和
\[
F(x) = f(x) + g(x)
\]
在点 $x = x_0$ 没有导数？
\end{prob}

\begin{prob}[M]{1015}
若 (a) 函数 $f(x)$ 在点 $x_0$ 有导数，而函数 $g(x)$ 在此点没有导数；(b) 函数 $f(x)$ 和 $g(x)$ 二者在点 $x_0$ 都没有导数，可否断定它们的积
\[
F(x) = f(x)g(x)
\]
在点 $x = x_0$ 没有导数？

研究两个例子：
(a) $f(x) = x$, $g(x) = |x|$, $x_0 = 0$； \qquad (b) $f(x) = |x|$, $g(x) = |x|$, $x_0 = 0$．
\end{prob}

\begin{prob}[L]{1016}
若 (a) 函数 $f(x)$ 于点 $x = g(x_0)$ 有导数，而函数 $g(x)$ 于点 $x = x_0$ 没有导数；(b) 函数 $f(x)$ 于点 $x = g(x_0)$ 没有导数，而函数 $g(x)$ 于点 $x = x_0$ 有导数；(c) 函数 $f(x)$ 于点 $x = g(x_0)$ 没有导数，函数 $g(x)$ 于点 $x = x_0$ 没有导数，则函数
\[
F(x) = f(g(x))
\]
于已知点 $x = x_0$ 的可微性怎样？

研究三个例子：
\begin{enumerate}[label=(\alph*), itemsep=2pt, topsep=3pt]
  \item $f(x) = x^2$, $g(x) = |x|$, $x_0 = 0$；
  \item $f(x) = |x|$, $g(x) = x^2$, $x_0 = 0$；
  \item $f(x) = 2x + |x|$, $g(x) = \dfrac{2}{3}x - \dfrac{1}{3}|x|$, $x_0 = 0$．
\end{enumerate}
\end{prob}

\begin{prob}[M]{1017}
函数
\[
y = x + \sqrt[3]{\sin x}
\]
的图像在哪些点处有竖直切线？作出此图像．
\end{prob}

\begin{prob}[M]{1018}
函数 $f(x)$ 在其间断点可否有：(a) 有限的导数；(b) 无穷的导数？研究例子：
\[
f(x) = \operatorname{sgn} x.
\]
\end{prob}

\begin{prob}[M]{1019}
若函数 $f(x)$ 在有限的区间 $(a, b)$ 上可微，且
\[
\lim_{x \to a} f(x) = \infty,
\]
则是否必有
\[
\text{(1) } \lim_{x \to a} f'(x) = \infty; \qquad \text{(2) } \overline{\lim}_{x \to a} |f'(x)| = +\infty\text{？}
\]
研究例子：$f(x) = \dfrac{1}{x} + \cos \dfrac{1}{x}$ 当 $x \to 0$ 时．
\end{prob}

\begin{prob}[M]{1020}
若函数 $f(x)$ 在有限的区间 $(a, b)$ 上可微，且
\[
\lim_{x \to a} f'(x) = \infty,
\]
则是否必有
\[
\lim_{x \to a} f(x) = \infty\text{？}
\]
研究例子：$f(x) = \sqrt[3]{x}$ 当 $x \to 0$ 时．
\end{prob}

\begin{prob}[M]{1021}
设函数 $f(x)$ 在区间 $(x_0, +\infty)$ 上可微，且 $\lim_{x \to +\infty} f(x)$ 存在．由此能否推出 $\lim_{x \to +\infty} f'(x)$ 存在？研究例子：
\[
f(x) = \frac{\sin(x^2)}{x}.
\]
\end{prob}

\begin{prob}[M]{1022}
设有界函数 $f(x)$ 在区间 $(x_0, +\infty)$ 上可微，且 $\lim_{x \to +\infty} f'(x)$ 存在．由此可否推出有限的或无穷的 $\lim_{x \to +\infty} f(x)$ 存在？研究例子：
\[
f(x) = \cos(\ln x).
\]
\end{prob}

\begin{prob}[S]{1023}
对函数之间的不等式可否逐项微分？
\end{prob}

\begin{prob}[L]{1024}
导出求和公式：
\begin{enumerate}[label=(\alph*), itemsep=2pt, topsep=3pt]
  \item $P_n = 1 + 2x + 3x^2 + \cdots + nx^{n-1}$，\\
  $Q_n = 1^2 + 2^2 x + 3^2 x^2 + \cdots + n^2 x^{n-1}$；
  \probhint{研究 $(x + x^2 + \cdots + x^n)'$．}
  \item $S_n = \sin x + \sin 2x + \cdots + \sin nx$，\\
  $T_n = \cos x + 2\cos 2x + \cdots + n\cos nx$．
\end{enumerate}
\end{prob}

\begin{prob}[M]{1025}
导出求和公式：
\[
S_n = \cosh x + 2\cosh 2x + \cdots + n\cosh nx.
\]
\probhint{$S_n = (\sinh x + \sinh 2x + \cdots + \sinh nx)'$．}
\end{prob}

\begin{prob}[M]{1026}
利用恒等式
\[
\cos \frac{x}{2} \cos \frac{x}{4} \cdots \cos \frac{x}{2^n} = \frac{\sin x}{2^n \sin \frac{x}{2^n}}
\]
推出求和公式：
\[
S_n = \frac{1}{2} \tan \frac{x}{2} + \frac{1}{4} \tan \frac{x}{4} + \cdots + \frac{1}{2^n} \tan \frac{x}{2^n}.
\]
\end{prob}

\begin{prob}[M]{1027}
证明：可微偶函数的导数为奇函数，而可微奇函数的导数为偶函数．给出这个事实的几何解释．
\end{prob}

\begin{prob}[M]{1028}
证明：可微周期函数的导数仍为具有相同周期的周期函数．
\end{prob}

\begin{prob}[S]{1029}
若圆半径以 $2\,\mathrm{cm/s}$ 的速度匀速增加，则当圆半径 $R = 10\,\mathrm{cm}$ 时，圆面积增加的速度如何？
\end{prob}

\begin{prob}[M]{1030}
若矩形的一边以 $1\,\mathrm{m/s}$ 的速度缩短，另一边以 $2\,\mathrm{m/s}$ 的速度伸长，则当这两条边达到 $x = 20\,\mathrm{m}$ 和 $y = 15\,\mathrm{m}$ 时，矩形的面积和对角线变化的速度如何？
\end{prob}

\begin{prob}[M]{1031}
二轮船 $A$ 和 $B$ 从同一码头同时出发，$A$ 船往北，$B$ 船往东．若 $A$ 船的速度为 $30\,\mathrm{km/h}$，$B$ 船的速度为 $40\,\mathrm{km/h}$，问二船间的距离增加的速度如何？
\end{prob}

\begin{prob}[L]{1032}
设
\[
f(x) = \begin{cases}
x, & 0 \leqslant x \leqslant 2, \\
2x - 2, & 2 < x < +\infty,
\end{cases}
\]
又设 $S(x)$ 表示由曲线 $y = f(x)$，轴 $Ox$ 及通过点 $x$ ($x \geqslant 0$) 且垂直于 $Ox$ 的直线三者围成的面积．写出函数 $S(x)$ 的解析表达式，求出导数 $S'(x)$，并作出函数 $y = S'(x)$ 的图像．
\end{prob}

\begin{prob}[L]{1033}
函数 $S(x)$ 是由圆弧 $y = \sqrt{a^2 - x^2}$、轴 $Ox$ 及分别通过点 $O$ 和 $x$ ($|x| \leqslant a$) 且垂直于轴 $Ox$ 的两条直线四者围成的面积．写出函数 $S(x)$ 的解析表达式，求出导数 $S'(x)$，并作此导数 $y = S'(x)$ 的图像．
\end{prob}

\probsec{§2. 反函数的导数. 用参数形式给出的函数的导数. 隐函数的导数}

\begin{prob}[M]{1034}
证明：由方程 $y^3 + 3y = x$ 确定的单值函数 $y = y(x)$ 存在，并求它的导数 $y'_x$．
\end{prob}

\begin{prob}[M]{1035}
证明：由方程 $y - \varepsilon \sin y = x$ ($0 \leqslant \varepsilon < 1$) 确定的单值函数 $y = y(x)$ 存在，并求它的导数 $y'_x$．
\end{prob}

\begin{prob}[L]{1036}
设：
\begin{enumerate}[label=(\alph*), itemsep=2pt, topsep=3pt]
  \item $y = x + \ln x$ ($x > 0$)；
  \item $y = x + \mathrm{e}^x$；
  \item $y = \sinh x$；
  \item $y = \tanh x$．
\end{enumerate}
求反函数 $x = x(y)$ 的存在域，并求它们的导数．
\end{prob}

\begin{prob}[L]{1037}
设：
\begin{enumerate}[label=(\alph*), itemsep=2pt, topsep=3pt]
  \item $y = 2x^2 - x^4$；
  \item $y = \dfrac{x^2}{1 + x^2}$；
  \item $y = 2\mathrm{e}^{-x} - \mathrm{e}^{-2x}$．
\end{enumerate}
选出反函数 $x = x(y)$ 的单值连续的各支，求它们的导数并作其图像．
\end{prob}

\begin{prob}[L]{1038}
作出函数 $y = y(x)$ 的略图，并求其导数 $y'_x(x)$．设：$x = -1 + 2t - t^2$，$y = 2 - 3t + t^3$．当 $x = 0$ 及 $x = -1$ 时，$y'_x(x)$ 等于什么？在何点 $M(x, y)$ 的导数 $y'_x(x) = 0$？
\end{prob}

\probgroup{求导数 $y'_x$（参数是正数），设：}

\begin{prob}[S]{1039}
$x = \sqrt[3]{1 - \sqrt{t}}$，$y = \sqrt{1 - \sqrt[3]{t}}$．
\end{prob}

\begin{prob}[S]{1040}
$x = \sin^2 t$，$y = \cos^2 t$．
\end{prob}

\begin{prob}[S]{1041}
$x = a\cos t$，$y = b\sin t$．
\end{prob}

\begin{prob}[S]{1042}
$x = a\cosh t$，$y = b\sinh t$．
\end{prob}

\begin{prob}[S]{1043}
$x = a\cos^3 t$，$y = a\sin^3 t$．
\end{prob}

\begin{prob}[S]{1044}
$x = a(t - \sin t)$，$y = a(1 - \cos t)$．
\end{prob}

\begin{prob}[S]{1045}
$x = \mathrm{e}^{2t} \cos^2 t$，$y = \mathrm{e}^{2t} \sin^2 t$．
\end{prob}

\begin{prob}[S]{1046}
$x = \arcsin \dfrac{t}{\sqrt{1 + t^2}}$，$y = \arccos \dfrac{1}{\sqrt{1 + t^2}}$．
\end{prob}

\begin{prob}[M]{1047}
证明：由方程组
\[
\begin{cases}
x = 2t + |t|, \\
y = 5t^2 + 4t|t|
\end{cases}
\]
所确定的函数 $y = y(x)$ 当 $t = 0$ 时可微，但它的导数在此点不能用普通的公式求得．
\end{prob}

\probgroup{求下列隐函数的导数 $y'_x$：}

\begin{prob}[M]{1048}
$x^2 + 2xy - y^2 = 2x$．

当 $x = 2$ 与 $y = 4$ 时及当 $x = 2$ 与 $y = 0$ 时，$y'$ 等于什么？
\end{prob}

\begin{prob}[S]{1049}
$y^2 = 2px$（抛物线）．
\end{prob}

\begin{prob}[S]{1050}
$\dfrac{x^2}{a^2} + \dfrac{y^2}{b^2} = 1$（椭圆）．
\end{prob}

\begin{prob}[S]{1051}
$\sqrt{x} + \sqrt{y} = \sqrt{a}$（抛物线）．
\end{prob}

\begin{prob}[S]{1052}
$x^{\frac{2}{3}} + y^{\frac{2}{3}} = a^{\frac{2}{3}}$（星形线）．
\end{prob}

\begin{prob}[S]{1053}
$\arctan \dfrac{y}{x} = \ln \sqrt{x^2 + y^2}$（对数螺线）．
\end{prob}

\begin{prob}[L]{1054}
设 $y'_x$，设：
\begin{enumerate}[label=(\alph*), itemsep=2pt, topsep=3pt]
  \item $r = a\varphi$（阿基米德螺线）；
  \item $r = a(1 + \cos\varphi)$（心脏线）；
  \item $r = a\mathrm{e}^{m\varphi}$（对数螺线），
\end{enumerate}
其中 $r = \sqrt{x^2 + y^2}$ 及 $\varphi = \arctan \dfrac{y}{x}$ 是极坐标．
\end{prob}

\probsec{§3. 导数的几何意义}

\begin{prob}[L]{1055}
写出曲线
\[
y = (x + 1)\sqrt[3]{3 - x}
\]
在下列各点处的切线和法线方程：
\begin{enumerate}[label=(\alph*), itemsep=2pt, topsep=3pt]
  \item $A(-1, 0)$；
  \item $B(2, 3)$；
  \item $C(3, 0)$．
\end{enumerate}
\end{prob}

\begin{prob}[M]{1056}
在曲线
\[
y = 2 + x - x^2
\]
上的哪些点，其切线 (a) 平行于 $Ox$ 轴；(b) 平行于第一象限角的平分线？
\end{prob}

\begin{prob}[M]{1057}
证明：抛物线
\[
y = a(x - x_1)(x - x_2) \quad (a \neq 0, x_1 < x_2)
\]
与 $Ox$ 轴相交所成的两角 $\alpha$ 及 $\beta$ $\left(0 < \alpha < \dfrac{\pi}{2}, 0 < \beta < \dfrac{\pi}{2}\right)$ 彼此相等．
\end{prob}

\begin{prob}[M]{1058}
在曲线
\[
y = 2\sin x \quad (-\pi \leqslant x \leqslant \pi)
\]
上求出“曲线的坡度”（即 $|y'|$）大于 $1$ 的部分.
\end{prob}

\begin{prob}[M]{1059}
函数
\[
y = x \quad \text{及} \quad y_1 = x + 0.01\sin 1000\pi x
\]
二者相差不大于 $0.01$. 这些函数的导数的差的最大值如何？作出相应的图像.
\end{prob}

\begin{prob}[S]{1060}
曲线 $y = \ln x$ 与 $Ox$ 轴相交的角如何？
\end{prob}

\begin{prob}[S]{1061}
曲线 $y = x^2$ 与 $x = y^2$ 相交的角如何？
\end{prob}

\begin{prob}[S]{1062}
曲线 $y = \sin x$ 与 $y = \cos x$ 相交的角如何？
\end{prob}

\begin{prob}[M]{1063.1}
如何选择参数 $n$，可使曲线
\[
y = \arctan nx \quad (n > 0)
\]
与 $Ox$ 轴相交所成的角大于 $89^\circ$？
\end{prob}

\begin{prob}[M]{1063.2}
证明：曲线 $y = |x|^\alpha$ 在不同情形下与下列坐标轴相切：
(a) 当 $0 < \alpha < 1$ 时与 $Oy$ 轴相切；\qquad (b) 当 $1 < \alpha < +\infty$ 时与 $Ox$ 轴相切.
\end{prob}

\begin{prob}[M]{1063.3}
证明：对于函数
\[
y = \begin{cases}
|x|^\alpha, & x \neq 0, \, \alpha \neq 0, \\
1, & x = 0
\end{cases}
\]
的图像，经过点 $A(0, 1)$ 的割线的极限位置是 $Oy$ 轴.
\end{prob}

\begin{prob}[M]{1064}
求出曲线：(a) $y = \sqrt{1 - \mathrm{e}^{-a^2 x^2}}$ 于点 $x = 0$ 处；(b) $y = \arcsin \dfrac{2x}{1 + x^2}$ 于点 $x = 1$ 处的左切线与右切线间的夹角.
\end{prob}

\begin{prob}[M]{1065}
证明：对数螺线 $r = a\mathrm{e}^{m\varphi}$（$a$ 及 $m$ 为常数）的切线与切点的径向量所成的角度为一常量.
\end{prob}

\begin{prob}[M]{1066}
求曲线 $y = ax^n$ 的次切线长，由此给出作此曲线的切线的方法.
\end{prob}

\begin{prob}[M]{1067}
证明：对于抛物线 $y^2 = 2px$，(a) 次切线长等于切点的横坐标（绝对值）的两倍；(b) 次法线为一常量. 给出作抛物线的切线的方法.
\end{prob}

\begin{prob}[M]{1068}
证明：指数曲线 $y = a^x$ ($a > 0$) 有定长的次切线. 给出作指数曲线的切线的方法.
\end{prob}

\begin{prob}[M]{1069}
求悬链线
\[
y = a\cosh \frac{x}{a}
\]
上任何点 $M(x_0, y_0)$ 处的法线长.
\end{prob}

\begin{prob}[M]{1070}
证明：星形线
\[
x^{\frac{2}{3}} + y^{\frac{2}{3}} = a^{\frac{2}{3}} \quad (a > 0)
\]
的切线介于坐标轴间的部分的长为一常量.
\end{prob}

\begin{prob}[S]{1071}
若抛物线 $y = ax^2 + bx + c$ 与 $Ox$ 轴相切，则系数 $a, b, c$ 间的关系如何？
\end{prob}

\begin{prob}[M]{1072}
在什么条件下立方抛物线
\[
y = x^3 + px + q
\]
与 $Ox$ 轴相切？
\end{prob}

\begin{prob}[S]{1073}
当参数 $a$ 为何值时，抛物线 $y = ax^2$ 与曲线 $y = \ln x$ 相切？
\end{prob}

\begin{prob}[M]{1074}
证明：曲线
\[
y = f(x) \quad (f(x) > 0)
\]
和
\[
y = f(x)\sin ax
\]
于公共点彼此相切，其中 $f(x)$ 为可微函数.
\end{prob}

\begin{prob}[M]{1075}
证明：双曲线族 $x^2 - y^2 = a$ 及 $xy = b$ 形成一正交网，即这两族曲线成直角相交.
\end{prob}

\begin{prob}[M]{1076}
证明：抛物线族
\[
y^2 = 4a(a - x) \quad (a > 0)
\]
及
\[
y^2 = 4b(b + x) \quad (b > 0)
\]
形成正交网.
\end{prob}

\begin{prob}[M]{1077}
写出曲线 $x = 2t - t^2$ 及 $y = 3t - t^3$ 在下列各点的切线和法线的方程：
(a) $t = 0$；\qquad (b) $t = 1$.
\end{prob}

\begin{prob}[L]{1078}
写出曲线
\[
x = \frac{2t + t^2}{1 + t^3}, \quad y = \frac{2t - t^2}{1 + t^3}
\]
在下列各点的切线与法线方程：(a) $t = 0$，(b) $t = 1$，(c) $t = \infty$.
\end{prob}

\begin{prob}[M]{1079}
写出摆线（旋轮线）
\[
x = a(t - \sin t), \quad y = a(1 - \cos t)
\]
上任意一点 $t = t_0$ 处的切线方程. 给出摆线的切线的作法.
\end{prob}

\begin{prob}[M]{1080}
证明：曳物线
\[
x = a\left(\ln \tan \frac{t}{2} + \cos t\right), \quad y = a\sin t \quad (a > 0, \, 0 < t < \pi)
\]
有长度不变的切线.
\end{prob}

\probgroup{写出下列曲线在指定点的切线与法线方程：}

\begin{prob}[S]{1081}
$\dfrac{x^2}{100} + \dfrac{y^2}{64} = 1, \; M(6, 6.4)$.
\end{prob}

\begin{prob}[S]{1082}
$xy + \ln y = 1, \; M(1, 1)$.
\end{prob}

\probsec{§4. 函数的微分}

\begin{prob}[L]{1083}
设
(a) $\Delta x = 1$；\qquad (b) $\Delta x = 0.1$；\qquad (c) $\Delta x = 0.01$.
对于函数
\[
f(x) = x^3 - 2x + 1,
\]
求：(1) $\Delta f(1)$，(2) $\mathrm{d}f(1)$，并进行对比.
\end{prob}

\begin{prob}[L]{1084}
设运动方程由公式
\[
x = 5t^2
\]
给出，其中 $t$ 以 $\mathrm{s}$（秒）来度量，$x$ 以 $\mathrm{m}$（米）来度量. 若 (a) $\Delta t = 1\,\mathrm{s}$，(b) $\Delta t = 0.1\,\mathrm{s}$，(c) $\Delta t = 0.001\,\mathrm{s}$，对于 $t = 2\,\mathrm{s}$ 的时刻，求路程的增量 $\Delta x$ 及微分 $\mathrm{d}x$，并作比较.
\end{prob}

\probgroup{求下列函数 $y$ 的微分：}

\begin{prob}[S]{1085}
$y = \dfrac{1}{x}$.
\end{prob}

\begin{prob}[S]{1086}
$y = \dfrac{1}{a}\arctan \dfrac{x}{a} \; (a \neq 0)$.
\end{prob}

\begin{prob}[S]{1087}
$y = \dfrac{1}{2a}\ln\left|\dfrac{x - a}{x + a}\right|$.
\end{prob}

\begin{prob}[S]{1088}
$y = \ln\left|x + \sqrt{x^2 + a}\right|$.
\end{prob}

\begin{prob}[S]{1089}
$y = \arcsin \dfrac{x}{a} \; (a \neq 0)$.
\end{prob}

\begin{prob}[L]{1090}
求：
\begin{enumerate}[label=(\alph*), itemsep=2pt, topsep=3pt]
\item $\mathrm{d}(x\mathrm{e}^x)$；
\item $\mathrm{d}(\sin x - x\cos x)$；
\item $\mathrm{d}\left(\dfrac{1}{x^3}\right)$；
\item $\mathrm{d}\left(\dfrac{\ln x}{\sqrt{x}}\right)$；
\item $\mathrm{d}(\sqrt{a^2 + x^2})$；
\item $\mathrm{d}\left(\dfrac{x}{\sqrt{1 - x^2}}\right)$；
\item $\mathrm{d}\ln(1 - x^2)$；
\item $\mathrm{d}\left(\arccos \dfrac{1}{|x|}\right)$；
\item $\mathrm{d}\left[\dfrac{\sin x}{2\cos^2 x} + \dfrac{1}{2}\ln\left|\tan\left(\dfrac{x}{2} + \dfrac{\pi}{4}\right)\right|\right]$.
\end{enumerate}
\end{prob}

\probgroup{设 $u, v, w$ 为 $x$ 的可微函数，求下列函数 $y$ 的微分：}

\begin{prob}[S]{1091}
$y = uvw$.
\end{prob}

\begin{prob}[S]{1092}
$y = \dfrac{u}{v^2}$.
\end{prob}

\begin{prob}[S]{1093}
$y = \dfrac{1}{\sqrt{u^2 + v^2}}$.
\end{prob}

\begin{prob}[S]{1094}
$y = \arctan \dfrac{u}{v}$.
\end{prob}

\begin{prob}[S]{1095}
$y = \ln\sqrt{u^2 + v^2}$.
\end{prob}

\begin{prob}[L]{1096}
求：
\begin{enumerate}[label=(\alph*), itemsep=2pt, topsep=3pt]
\item $\dfrac{\mathrm{d}}{\mathrm{d}x^3}(x^3 - 2x^6 - x^9)$；
\item $\dfrac{\mathrm{d}}{\mathrm{d}x^2}\left(\dfrac{\sin x}{x}\right)$；
\item $\dfrac{\mathrm{d}(\sin x)}{\mathrm{d}(\cos x)}$；
\item $\dfrac{\mathrm{d}(\tan x)}{\mathrm{d}(\cot x)}$；
\item $\dfrac{\mathrm{d}(\arcsin x)}{\mathrm{d}(\arccos x)}$.
\end{enumerate}
\end{prob}

\begin{prob}[M]{1097}
有半径为 $R = 100\,\mathrm{cm}$ 及圆心角 $\alpha = 60^\circ$ 的扇形. 若 (a) 其半径 $R$ 增加 $1\,\mathrm{cm}$；(b) 角 $\alpha$ 减少 $30'$，则扇形面积的变化如何？求出精确的和近似的解.
\end{prob}

\begin{prob}[M]{1098}
摆的振动周期按下式确定：
\[
T = 2\pi\sqrt{\frac{l}{g}},
\]
其中 $l$ 为摆长，$g = 9.8\,\mathrm{m/s^2}$ 为重力加速度.
为了使周期 $T$ 增大 $0.05\,\mathrm{s}$，对摆长 $l = 20\,\mathrm{cm}$ 需要作多少修改？
\end{prob}

\probgroup{利用函数之微分代替其增量，求下列各式之近似值：}

\begin{prob}[S]{1099}
$\sqrt[3]{1.02}$.
\end{prob}

\begin{prob}[S]{1100}
$\sin 29^\circ$.
\end{prob}

\begin{prob}[S]{1101}
$\cos 151^\circ$.
\end{prob}

\begin{prob}[S]{1102}
$\arctan 1.05$.
\end{prob}

\begin{prob}[S]{1103}
$\lg 11$.
\end{prob}

\begin{prob}[L]{1104.1}
证明近似公式：
\[
\sqrt{a^2 + x} \approx a + \frac{x}{2a} \quad (a > 0),
\]
其中 $|x| \ll a$（正数 $A$ 和 $B$ 间的关系式 $A \ll B$ 表示 $A$ 远小于 $B$）. 利用此公式近似地计算：
(a) $\sqrt{5}$；\qquad (b) $\sqrt{34}$；\qquad (c) $\sqrt{120}$，
并与平方根表中的数值作比较.
\end{prob}

\begin{prob}[M]{1104.2}
证明公式：
\[
\sqrt{a^2 + x} = a + \frac{x}{2a} - r \quad (a > 0, \, x > 0),
\]
其中 $0 < r < \dfrac{x^2}{8a^3}$.
\end{prob}

\begin{prob}[L]{1105}
证明近似公式：
\[
\sqrt[n]{a^n + x} \approx a + \frac{x}{na^{n-1}} \quad (a > 0),
\]
其中 $|x| \ll a$. 利用此公式近似地计算：
(a) $\sqrt[3]{9}$；\qquad (b) $\sqrt[4]{80}$；\qquad (c) $\sqrt[7]{100}$；\qquad (d) $\sqrt[10]{1000}$.
\end{prob}

\begin{prob}[M]{1106}
正方形的边 $x = 2.4\,\mathrm{m} \pm 0.05\,\mathrm{m}$. 由此计算所得正方形的面积的相对误差和绝对误差如何？
\end{prob}

\begin{prob}[M]{1107}
为了在 $1\%$ 的精度下计算出球的体积，问度量球半径 $R$ 时所允许的相对误差如何？
\end{prob}

\begin{prob}[M]{1108}
为了确定重力加速度，可以利用摆的振动公式 $g = \dfrac{4\pi^2 l}{T^2}$，其中 $l$ 为摆长，$T$ 为振动周期. 当测量 (a) 长度 $l$，(b) 周期 $T$ 时，相对误差 $\delta$ 如何影响 $g$ 的值？
\end{prob}

\begin{prob}[M]{1109}
求数 $x$ ($x > 0$) 的常用对数（以 $10$ 为底）的绝对误差，设此数的相对误差等于 $\delta$.
\end{prob}

\begin{prob}[M]{1110}
证明：根据正切对数表求得的角度比根据具有同样多位小数的正弦对数表求得的角度更为精确.
\end{prob}

\probsec{§5. 高阶的导数和微分}

\probgroup{求 $y''$，设：}

\begin{prob}[S]{1111}
$y = x\sqrt{1 + x^2}$.
\end{prob}

\begin{prob}[S]{1112}
$y = \dfrac{x}{\sqrt{1 - x^2}}$.
\end{prob}

\begin{prob}[S]{1113}
$y = \mathrm{e}^{-x^2}$.
\end{prob}

\begin{prob}[S]{1114}
$y = \tan x$.
\end{prob}

\begin{prob}[S]{1115}
$y = (1 + x^2)\arctan x$.
\end{prob}

\begin{prob}[S]{1116}
$y = \frac{\arcsin x}{\sqrt{1 - x^2}}$.
\end{prob}

\begin{prob}[S]{1117}
$y = x\ln x$.
\end{prob}

\begin{prob}[S]{1118}
$y = \ln f(x)$.
\end{prob}

\begin{prob}[S]{1119}
$y = x[\sin(\ln x) + \cos(\ln x)]$.
\end{prob}

\begin{prob}[M]{1120}
设 $y = \mathrm{e}^{\sin x}\cos(\sin x)$, 求 $y(0)$, $y'(0)$ 及 $y''(0)$.
\end{prob}

\probgroup{设 $u = \varphi(x)$ 及 $v = \psi(x)$ 为二阶可微函数. 求 $y''$, 设:}

\begin{prob}[S]{1121}
$y = u^2$.
\end{prob}

\begin{prob}[S]{1122}
$y = \ln\frac{u}{v}$.
\end{prob}

\begin{prob}[S]{1123}
$y = \sqrt{u^2 + v^2}$.
\end{prob}

\begin{prob}[S]{1124}
$y = u^v$ ($u > 0$).
\end{prob}

\probgroup{设 $f(x)$ 为三阶可微函数. 求 $y''$ 及 $y'''$, 设:}

\begin{prob}[S]{1125}
$y = f(x^2)$.
\end{prob}

\begin{prob}[S]{1126}
$y = f\left(\frac{1}{x}\right)$.
\end{prob}

\begin{prob}[S]{1127}
$y = f(\mathrm{e}^x)$.
\end{prob}

\begin{prob}[S]{1128}
$y = f(\ln x)$.
\end{prob}

\begin{prob}[S]{1129}
$y = f(\varphi(x))$, 其中 $\varphi(x)$ 是充分多次可微函数.
\end{prob}

\begin{prob}[M]{1130}
对于以下两种情形: (a) $x$ 为自变量, (b) $x$ 为中间变量, 求函数 $y = \mathrm{e}^x$ 的 $\mathrm{d}^2 y$.
\end{prob}

\probgroup{若 $x$ 为自变量, 求 $\mathrm{d}^2 y$, 设:}

\begin{prob}[S]{1131}
$y = \sqrt{1 + x^2}$.
\end{prob}

\begin{prob}[S]{1132}
$y = \frac{\ln x}{x}$.
\end{prob}

\begin{prob}[S]{1133}
$y = x^x$.
\end{prob}

\probgroup{令 $u$ 及 $v$ 为变量 $x$ 的二阶可微函数, 求 $\mathrm{d}^2 y$, 设:}

\begin{prob}[S]{1134}
$y = uv$.
\end{prob}

\begin{prob}[S]{1135}
$y = \frac{u}{v}$.
\end{prob}

\begin{prob}[S]{1136}
$y = u^m v^n$ ($m$ 及 $n$ 为常数).
\end{prob}

\begin{prob}[S]{1137}
$y = a^u$ ($a > 0$).
\end{prob}

\begin{prob}[S]{1138}
$y = \ln\sqrt{u^2 + v^2}$.
\end{prob}

\begin{prob}[S]{1139}
$y = \arctan\frac{u}{v}$.
\end{prob}

\probgroup{对以参数形式给出的函数 $y = y(x)$ 的导数 $y'_x$, $y''_{x^2}$, $y'''_{x^3}$, 设:}

\begin{prob}[S]{1140}
$x = 2t - t^2$, $y = 3t - t^3$.
\end{prob}

\begin{prob}[S]{1141}
$x = a\cos t$, $y = a\sin t$.
\end{prob}

\begin{prob}[S]{1142}
$x = a(t - \sin t)$, $y = a(1 - \cos t)$.
\end{prob}

\begin{prob}[S]{1143}
$x = \mathrm{e}^t \cos t$, $y = \mathrm{e}^t \sin t$.
\end{prob}

\begin{prob}[S]{1144}
$x = f'(t)$, $y = t f'(t) - f(t)$.
\end{prob}

\begin{prob}[M]{1145}
设函数 $y = f(x)$ 充分多次可微. 求反函数 $x = f^{-1}(y)$ 的导数 $x'$, $x''$, $x'''$, $x^{(4)}$ (设这些导数都存在).
\end{prob}

\probgroup{求由下列隐函数给出的 $y = y(x)$ 的 $y'_x$, $y''_{x^2}$ 及 $y'''_{x^3}$:}

\begin{prob}[M]{1146}
$x^2 + y^2 = 25$. 在点 $M(3, 4)$ 的 $y'$, $y''$ 及 $y'''$ 等于什么?
\end{prob}

\begin{prob}[S]{1147}
$y^2 = 2px$.
\end{prob}

\begin{prob}[S]{1148}
$x^2 - xy + y^2 = 1$.
\end{prob}

\probgroup{求 $y'_x$ 及 $y''_{x^2}$, 设:}

\begin{prob}[S]{1149}
$y^2 + 2\ln y = x^4$.
\end{prob}

\begin{prob}[S]{1150}
$\sqrt{x^2 + y^2} = a\mathrm{e}^{\arctan\frac{y}{x}}$ ($a > 0$).
\end{prob}

\begin{prob}[M]{1151}
设函数 $f(x)$ 当 $x \leqslant x_0$ 时有定义且二阶可微. 应当如何选择系数 $a$, $b$ 及 $c$, 使函数
\[
F(x) = \begin{cases}
f(x), & x \leqslant x_0, \\
a(x - x_0)^2 + b(x - x_0) + c, & x > x_0
\end{cases}
\]
是二阶可微的函数?
\end{prob}

\begin{prob}[M]{1152}
质点作直线运动的规律为
\[
s = 10 + 20t - 5t^2.
\]
求其运动速度和加速度. 在 $t = 2$ 的时刻, 速度与加速度等于什么?
\end{prob}

\begin{prob}[M]{1153}
质点 $M(x, y)$ 沿圆周 $x^2 + y^2 = a^2$ 匀速运动, 运动一周的时间为 $T$. 若质点在时刻 $t = 0$ 位于点 $M_0(a, 0)$, 求质点的速度和加速度在 $Ox$ 轴上的投影 $v$ 和 $w$.
\end{prob}

\begin{prob}[L]{1154}
在重力场中, 质点 $M(x, y)$ 在竖直平面 $Oxy$ 内以初速度 $v_0$ 沿与水平面成 $\alpha$ 角的方向抛出. 建立运动方程 (忽略空气阻力), 并确定速度 $v$ 和加速度 $w$ 的大小及运动轨迹. 质点的最大上升高度和射程等于多少?
\end{prob}

\begin{prob}[M]{1155}
质点的运动方程为
\[
x = 4\sin\omega t - 3\cos\omega t, \quad y = 3\sin\omega t + 4\cos\omega t \quad (\omega \text{ 为常数}).
\]
求运动轨迹、速度与加速度的大小.
\end{prob}

\probgroup{求下列指定阶的导数:}

\begin{prob}[S]{1156}
$y = x(2x - 1)^2(x + 3)^3$, 求 $y^{(6)}$ 及 $y^{(7)}$.
\end{prob}

\begin{prob}[S]{1157}
$y = \frac{a}{x^m}$, 求 $y'''$.
\end{prob}

\begin{prob}[S]{1158}
$y = \sqrt{x}$, 求 $y^{(10)}$.
\end{prob}

\begin{prob}[S]{1159}
$y = \frac{x^2}{1 - x}$, 求 $y^{(8)}$.
\end{prob}

\begin{prob}[S]{1160}
$y = \frac{1 + x}{\sqrt{1 - x}}$, 求 $y^{(100)}$.
\end{prob}

\begin{prob}[S]{1161}
$y = x^2\mathrm{e}^{2x}$, 求 $y^{(20)}$.
\end{prob}

\begin{prob}[S]{1162}
$y = \frac{\mathrm{e}^x}{x}$, 求 $y^{(10)}$.
\end{prob}

\begin{prob}[S]{1163}
$y = x\ln x$, 求 $y^{(5)}$.
\end{prob}

\begin{prob}[S]{1164}
$y = \frac{\ln x}{x}$, 求 $y^{(5)}$.
\end{prob}

\begin{prob}[S]{1165}
$y = x^2\sin 2x$, 求 $y^{(50)}$.
\end{prob}

\begin{prob}[S]{1166}
$y = \frac{\cos 3x}{\sqrt[3]{1 - 3x}}$, 求 $y'''$.
\end{prob}

\begin{prob}[S]{1167}
$y = \sin x\sin 2x\sin 3x$, 求 $y^{(10)}$.
\end{prob}

\begin{prob}[S]{1168}
$y = x\sinh x$, 求 $y^{(100)}$.
\end{prob}

\begin{prob}[S]{1169}
$y = \mathrm{e}^x\cos x$, 求 $y^{(4)}$.
\end{prob}

\begin{prob}[S]{1170}
$y = \sin^2 x\ln x$, 求 $y^{(6)}$.
\end{prob}

\probgroup{在下列各例中, 视 $x$ 为自变量, 求指定阶的微分.}

\begin{prob}[S]{1171}
$y = x^5$, 求 $\mathrm{d}^5 y$.
\end{prob}

\begin{prob}[S]{1172}
$y = \frac{1}{\sqrt{x}}$, 求 $\mathrm{d}^3 y$.
\end{prob}

\begin{prob}[S]{1173}
$y = x\cos 2x$, 求 $\mathrm{d}^{10} y$.
\end{prob}

\begin{prob}[S]{1174}
$y = \mathrm{e}^x\ln x$, 求 $\mathrm{d}^4 y$.
\end{prob}

\begin{prob}[S]{1175}
$y = \cos x\cdot\cosh x$, 求 $\mathrm{d}^6 y$.
\end{prob}

\probgroup{设 $u$ 为 $x$ 的充分多次可微函数, 在下列各例中求指定阶的微分.}

\begin{prob}[S]{1176}
$y = u^2$, 求 $\mathrm{d}^{10} y$.
\end{prob}

\begin{prob}[S]{1177}
$y = \mathrm{e}^u$, 求 $\mathrm{d}^4 y$.
\end{prob}

\begin{prob}[S]{1178}
$y = \ln u$, 求 $\mathrm{d}^3 y$.
\end{prob}

\begin{prob}[M]{1179}
视 $x$ 为某个自变量的函数, 由函数 $y = f(x)$ 求 $\mathrm{d}^2 y$, $\mathrm{d}^3 y$ 及 $\mathrm{d}^4 y$.
\end{prob}

\begin{prob}[M]{1180}
以变量 $x$ 和 $y$ 的逐次微分来表示函数 $y = f(x)$ 的导数 $y''$ 及 $y'''$, 但不假定 $x$ 为自变量.
\end{prob}

\begin{prob}[M]{1181}
证明: 函数
\[
y = C_1\cos x + C_2\sin x,
\]
其中 $C_1$ 及 $C_2$ 为任意的常数, 满足方程
\[
y'' + y = 0.
\]
\end{prob}

\begin{prob}[M]{1182}
证明: 函数
\[
y = C_1\cosh x + C_2\sinh x,
\]
其中 $C_1$ 及 $C_2$ 为任意的常数, 满足方程
\[
y'' - y = 0.
\]
\end{prob}

\begin{prob}[M]{1183}
证明: 函数
\[
y = C_1\mathrm{e}^{\lambda_1 x} + C_2\mathrm{e}^{\lambda_2 x},
\]
其中 $C_1$ 及 $C_2$ 为任意的常数, $\lambda_1$ 及 $\lambda_2$ 为常数, 满足方程
\[
y'' - (\lambda_1 + \lambda_2)y' + \lambda_1\lambda_2 y = 0.
\]
\end{prob}

\begin{prob}[M]{1184}
证明: 函数
\[
y = x^n[C_1\cos(\ln x) + C_2\sin(\ln x)],
\]
其中 $C_1$ 及 $C_2$ 为任意的常数, $n$ 为常数, 满足方程
\[
x^2 y'' + (1 - 2n)xy' + (1 + n^2)y = 0.
\]
\end{prob}

\begin{prob}[L]{1185}
证明: 函数
\[
y = \mathrm{e}^{\frac{x}{\sqrt{2}}}\left(C_1\cos\frac{x}{\sqrt{2}} + C_2\sin\frac{x}{\sqrt{2}}\right) + \mathrm{e}^{-\frac{x}{\sqrt{2}}}\left(C_3\cos\frac{x}{\sqrt{2}} + C_4\sin\frac{x}{\sqrt{2}}\right),
\]
其中 $C_1, C_2, C_3$ 及 $C_4$ 为任意常数, 满足方程
\[
y^{(4)} + y = 0.
\]
\end{prob}

\begin{prob}[M]{1186}
证明: 若函数 $f(x)$ 有 $n$ 阶导数, 则
\[
[f(ax + b)]^{(n)} = a^n f^{(n)}(ax + b).
\]
\end{prob}

\begin{prob}[M]{1187}
若
\[
P(x) = a_0 x^n + a_1 x^{n-1} + \cdots + a_n,
\]
求 $P^{(n)}(x)$.
\end{prob}

\probgroup{求 $y^{(n)}$, 设:}

\begin{prob}[S]{1188}
$y = \frac{ax + b}{cx + d}$.
\end{prob}

\begin{prob}[S]{1189}
$y = \frac{1}{x(1 - x)}$.
\end{prob}

\begin{prob}[M]{1190}
$y = \frac{1}{x^2 - 3x + 2}$.
\probhint{分解函数为最简分式.}
\end{prob}

\begin{prob}[S]{1191}
$y = \frac{1}{\sqrt{1 - 2x}}$.
\end{prob}

\begin{prob}[S]{1192}
$y = \frac{x}{\sqrt[3]{1 + x}}$.
\end{prob}

\begin{prob}[S]{1193}
$y = \sin^2 x$.
\end{prob}

\begin{prob}[S]{1194}
$y = \cos^2 x$.
\end{prob}

\begin{prob}[S]{1195}
$y = \sin^3 x$.
\end{prob}

\begin{prob}[S]{1196}
$y = \cos^3 x$.
\end{prob}

\begin{prob}[S]{1197}
$y = \sin ax\sin bx$.
\end{prob}

\begin{prob}[S]{1198}
$y = \cos ax\cos bx$.
\end{prob}

\begin{prob}[S]{1199}
$y = \sin ax\cos bx$.
\end{prob}

\begin{prob}[S]{1200}
$y = \sin^2 ax\cos bx$.
\end{prob}

\begin{prob}[S]{1201}
$y = \sin^4 x + \cos^4 x$.
\end{prob}

\begin{prob}[S]{1202}
$y = x\cos ax$.
\end{prob}

\begin{prob}[S]{1203}
$y = x^2\sin ax$.
\end{prob}

\begin{prob}[S]{1204}
$y = (x^2 + 2x + 2)\mathrm{e}^{-x}$.
\end{prob}

\begin{prob}[S]{1205}
$y = \frac{\mathrm{e}^x}{x}$.
\end{prob}

\begin{prob}[S]{1206}
$y = \mathrm{e}^x\cos x$.
\end{prob}

\begin{prob}[S]{1207}
$y = \mathrm{e}^x\sin x$.
\end{prob}

\begin{prob}[S]{1208}
$y = \ln\frac{a + bx}{a - bx}$.
\end{prob}

\begin{prob}[S]{1209}
$y = \mathrm{e}^{ax} P(x)$, 其中 $P(x)$ 为多项式.
\end{prob}

\begin{prob}[S]{1210}
$y = x\sinh x$.
\end{prob}

\probgroup{求 $\mathrm{d}^n y$, 设:}

\begin{prob}[S]{1211}
$y = x^n\mathrm{e}^x$.
\end{prob}

\begin{prob}[S]{1212}
$y = \frac{\ln x}{x}$.
\end{prob}

\begin{prob}[L]{1213}
证明等式:
\begin{enumerate}[label=(\arabic*), itemsep=2pt, topsep=3pt]
  \item $[\mathrm{e}^{ax}\sin(bx + c)]^{(n)} = \mathrm{e}^{ax}(a^2 + b^2)^{\frac{n}{2}}\sin(bx + c + n\varphi)$,
  \item $[\mathrm{e}^{ax}\cos(bx + c)]^{(n)} = \mathrm{e}^{ax}(a^2 + b^2)^{\frac{n}{2}}\cos(bx + c + n\varphi)$,
\end{enumerate}
其中 $\sin\varphi = \frac{b}{\sqrt{a^2 + b^2}}$, $\cos\varphi = \frac{a}{\sqrt{a^2 + b^2}}$.
\end{prob}

\begin{prob}[M]{1214}
求 $y^{(n)}$, 设:
(a) $y = \cosh ax\cos bx$; \qquad (b) $y = \cosh ax\sin bx$.
\end{prob}

\begin{prob}[L]{1215}
将函数
\[
f(x) = \sin^{2p} x,
\]
其中 $p$ 为正整数, 化为三角多项式:
\[
f(x) = \sum_{k=0}^{p} A_k\cos 2kx,
\]
以求 $f^{(n)}(x)$.
\probhint{令 $\sin x = \frac{1}{2\mathrm{i}}(t - \bar{t})$, 其中 $t = \cos x + \mathrm{i}\sin x$ 及 $\bar{t} = \cos x - \mathrm{i}\sin x$, 并且利用棣莫弗公式来计算.}
\end{prob}

\begin{prob}[L]{1216}
设:
\begin{enumerate}[label=(\alph*), itemsep=2pt, topsep=3pt]
  \item $f(x) = \sin^{2p+1} x$;
  \item $f(x) = \cos^{2p} x$;
  \item $f(x) = \cos^{2p+1} x$,
\end{enumerate}
其中 $p$ 为正整数, 求 $f^{(n)}(x)$ (参阅前题).
\end{prob}

\probgroup{若
\[
f(x) = f_1(x) + \mathrm{i} f_2(x),
\]
其中 $\mathrm{i}$ 是虚数单位, $f_1(x)$ 和 $f_2(x)$ 为实变量 $x$ 的实函数, 则按定义有:
\[
f'(x) = f'_1(x) + \mathrm{i} f'_2(x).
\]}

\begin{prob}[L]{1217}
利用恒等式
\[
\frac{1}{x^2 + 1} = \frac{1}{2\mathrm{i}}\left(\frac{1}{x - \mathrm{i}} - \frac{1}{x + \mathrm{i}}\right)
\]
证明:
\[
\left(\frac{1}{x^2 + 1}\right)^{(n)} = \frac{(-1)^n n!}{(1 + x^2)^{\frac{n+1}{2}}}\sin[(n + 1)\operatorname{arccot} x].
\]
\probhint{利用棣莫弗公式.}
\end{prob}

\begin{prob}[M]{1218}
求函数 $f(x) = \arctan x$ 的 $n$ 阶导数.
\end{prob}

\probgroup{求 $f^{(n)}(0)$, 设:}

\begin{prob}[M]{1219}
(a) $f(x) = \frac{1}{(1 - 2x)(1 + x)}$; \qquad (b) $f(x) = \frac{x}{\sqrt{1 - x}}$.
\end{prob}

\begin{prob}[L]{1220}
\begin{enumerate}[label=(\alph*), itemsep=2pt, topsep=3pt]
  \item $f(x) = x^2\mathrm{e}^{ax}$;
  \item $f(x) = \arctan x$;
  \item $f(x) = \arcsin x$.
\end{enumerate}
\end{prob}

\begin{prob}[M]{1221}
(a) $f(x) = \cos(m\arcsin x)$; \qquad (b) $f(x) = \sin(m\arcsin x)$.
\end{prob}

\begin{prob}[M]{1222}
(a) $f(x) = (\arctan x)^2$; \qquad (b) $f(x) = (\arcsin x)^2$.
\end{prob}

\begin{prob}[M]{1223}
设
\[
f(x) = (x - a)^n\varphi(x),
\]
其中函数 $\varphi(x)$ 在点 $a$ 的邻域内有 $n - 1$ 阶的连续导数, 求 $f^{(n)}(a)$.
\end{prob}

\begin{prob}[M]{1224}
证明: 函数
\[
f(x) = \begin{cases}
x^{2n}\sin\dfrac{1}{x}, & x \neq 0, \\
0, & x = 0
\end{cases}
\]
($n$ 为正整数), 于点 $x = 0$ 有一直到 $n$ 阶的导数, 而无 $n + 1$ 阶导数.
\end{prob}

\begin{prob}[M]{1225}
证明: 函数
\[
f(x) = \begin{cases}
\mathrm{e}^{-1/x^2}, & x \neq 0, \\
0, & x = 0
\end{cases}
\]
在 $x = 0$ 处是无穷次可微的. 作出此函数的图像.
\end{prob}

\begin{prob}[M]{1226}
证明: 切比雪夫多项式
\[
T_m(x) = \frac{1}{2^{m-1}}\cos(m\arccos x) \quad (m = 1, 2, \cdots)
\]
满足方程
\[
(1 - x^2)T''_m(x) - xT'_m(x) + m^2 T_m(x) = 0.
\]
\end{prob}

\begin{prob}[L]{1227}
证明: 勒让德多项式
\[
P_m(x) = \frac{1}{2^m m!}[(x^2 - 1)^m]^{(m)} \quad (m = 0, 1, 2, \cdots)
\]
满足方程
\[
(1 - x^2)P''_m(x) - 2xP'_m(x) + m(m + 1)P_m(x) = 0.
\]
\probhint{把等式 $(x^2 - 1)u' = 2mxu$ 作 $m + 1$ 次微分, 其中 $u = (x^2 - 1)^m$.}
\end{prob}

\begin{prob}[L]{1228}
切比雪夫-拉盖尔多项式定义如下:
\[
L_m(x) = \mathrm{e}^x(x^m\mathrm{e}^{-x})^{(m)} \quad (m = 0, 1, 2, \cdots).
\]
求多项式 $L_m(x)$ 的显式表达式.
证明: $L_m(x)$ 满足方程
\[
xL''_m(x) + (1 - x)L'_m(x) + mL(x) = 0.
\]
\probhint{利用等式 $xu' + (x - m)u = 0$，其中 $u = x^m \mathrm{e}^{-x}$.}
\end{prob}

\begin{prob}[M]{1229}
设 $y = f(u)$ 及 $u = \varphi(x)$，其中 $f(u)$ 及 $\varphi(x)$ 为 $n$ 阶可微函数. 证明：
\[
\frac{\mathrm{d}^n y}{\mathrm{d}x^n} = \sum_{k=1}^{n} A_k(x) f^{(k)}(u),
\]
其中系数 $A_k(x)$ $(k = 0, 1, \cdots, n)$ 与函数 $f(u)$ 无关.
\end{prob}

\begin{prob}[M]{1230}
证明：对于复合函数 $y = f(x^2)$ 的 $n$ 阶导数，成立公式
\[
\frac{\mathrm{d}^n y}{\mathrm{d}x^n} = (2x)^n f^{(n)}(x^2) + \frac{n(n-1)}{1!}(2x)^{n-2} f^{(n-1)}(x^2) + \frac{n(n-1)(n-2)(n-3)}{2!}(2x)^{n-4} f^{(n-2)}(x^2) + \cdots.
\]
\end{prob}

\begin{prob}[L]{1231}
切比雪夫--埃尔米特多项式定义如下：
\[
H_m(x) = (-1)^m \mathrm{e}^{x^2} (\mathrm{e}^{-x^2})^{(m)} \quad (m = 0, 1, 2, \cdots).
\]
求多项式 $H_m(x)$ 的显式表达式.

证明：$H_m(x)$ 满足方程
\[
H_m''(x) - 2x H_m'(x) + 2m H_m(x) = 0.
\]
\probhint{利用等式 $u' + 2xu = 0$，其中 $u = \mathrm{e}^{-x^2}$.}
\end{prob}

\begin{prob}[M]{1232.1}
证明等式：
\[
\left(x^{n-1} \mathrm{e}^{\frac{1}{x}}\right)^{(n)} = \frac{(-1)^n}{x^{n+1}} \mathrm{e}^{\frac{1}{x}}.
\]
\probhint{运用数学归纳法.}
\end{prob}

\begin{prob}[M]{1232.2}
证明公式：
\begin{enumerate}[label=(\alph*), itemsep=2pt, topsep=3pt]
  \item $\dfrac{\mathrm{d}^n}{\mathrm{d}x^n}(x^n \ln x) = n!\left(\ln x + \sum\limits_{k=1}^{n} \dfrac{1}{k}\right) \quad (x > 0)$;
  \item $\dfrac{\mathrm{d}^{2n}}{\mathrm{d}x^{2n}}\left(\dfrac{\sin x}{x}\right) = \dfrac{(2n)!}{x^{2n+1}}[C_n(x)\sin x - S_n(x)\cos x]$，其中
  \[
  C_n(x) = 1 - \frac{x^2}{2!} + \cdots + (-1)^n \frac{x^{2n}}{(2n)!}, \quad S_n(x) = x - \frac{x^3}{3!} + \cdots + (-1)^{n-1} \frac{x^{2n-1}}{(2n-1)!}.
  \]
\end{enumerate}
\end{prob}

\begin{prob}[M]{1233}
设 $\dfrac{\mathrm{d}}{\mathrm{d}x} = D$ 表示微分算子，
\[
f(D) = \sum_{k=0}^{n} p_k(x) D^k
\]
为微分符号的多项式，其中 $p_k(x)$ $(k = 0, 1, \cdots, n)$ 为 $x$ 的某连续函数. 证明：
\[
f(D)\{\mathrm{e}^{\lambda x} u(x)\} = \mathrm{e}^{\lambda x} f(D + \lambda) u(x),
\]
其中 $\lambda$ 为常数.
\end{prob}

\begin{prob}[M]{1234}
证明：若在方程
\[
\sum_{k=0}^{n} a_k x^k y_x^{(k)} = 0
\]
中令
\[
x = \mathrm{e}^t,
\]
其中 $t$ 为自变量，则此方程化为
\[
\sum_{k=0}^{n} a_k D(D - 1)\cdots(D - k + 1) y = 0,
\]
其中 $D = \dfrac{\mathrm{d}}{\mathrm{d}t}$.
\end{prob}

\probsec{§6. 罗尔定理、拉格朗日定理和柯西定理}

\begin{prob}[S]{1235}
检验罗尔定理对于函数
\[
f(x) = (x - 1)(x - 2)(x - 3)
\]
的正确性.
\end{prob}

\begin{prob}[M]{1236}
函数 $f(x) = 1 - \sqrt[3]{x^2}$ 当 $x_1 = -1$ 及 $x_2 = 1$ 时为零，但是当 $-1 \leqslant x \leqslant 1$ 时，$f'(x) \neq 0$. 说明与罗尔定理表面上的矛盾.
\end{prob}

\begin{prob}[M]{1237}
设函数 $f(x)$ 在有限或无穷的区间 $(a, b)$ 中的任意一点有有限的导数 $f'(x)$，且
\[
\lim_{x \to a+0} f(x) = \lim_{x \to b-0} f(x).
\]
证明：
\[
f'(c) = 0,
\]
其中 $c$ 为区间 $(a, b)$ 中的某点.
\end{prob}

\begin{prob}[M]{1238}
设函数 $f(x)$：(1) 在闭区间 $[x_0, x_n]$ 上有定义且有 $n - 1$ 阶的连续导数 $f^{(n-1)}(x)$；(2) 在区间 $(x_0, x_n)$ 内有 $n$ 阶导数 $f^{(n)}(x)$；(3) 下列等式成立：
\[
f(x_0) = f(x_1) = \cdots = f(x_n) \quad (x_0 < x_1 < \cdots < x_n).
\]
证明：在区间 $(x_0, x_n)$ 内最少存在一点 $\xi$，使
\[
f^{(n)}(\xi) = 0.
\]
\end{prob}

\begin{prob}[M]{1239}
设函数 $f(x)$：(1) 在闭区间 $[a, b]$ 上有定义且有 $p + q$ 阶的连续导数 $f^{(p+q)}(x)$；(2) 在区间 $(a, b)$ 内有 $p + q + 1$ 阶的导数 $f^{(p+q+1)}(x)$；(3) 下列等式成立：
\[
f(a) = f'(a) = \cdots = f^{(p)}(a) = 0,
\]
\[
f(b) = f'(b) = \cdots = f^{(q)}(b) = 0.
\]
证明：在此种情形下
\[
f^{(p+q+1)}(c) = 0,
\]
其中 $c$ 为区间 $(a, b)$ 内的某点.
\end{prob}

\begin{prob}[M]{1240}
证明：若具实系数 $a_k$ $(k = 0, 1, \cdots, n)$ 的多项式
\[
P_n(x) = a_0 x^n + a_1 x^{n-1} + \cdots + a_n \quad (a_0 \neq 0)
\]
之一切根为实数，则其逐次的导数 $P_n'(x), P_n''(x), \cdots, P_n^{(n-1)}(x)$ 也仅有实根.
\end{prob}

\begin{prob}[M]{1241}
证明：勒让德多项式
\[
P_n(x) = \frac{1}{2^n n!} \frac{\mathrm{d}^n}{\mathrm{d}x^n}\{(x^2 - 1)^n\}
\]
的一切根都是实数且包含于区间 $(-1, 1)$ 中.
\end{prob}

\begin{prob}[M]{1242}
证明：切比雪夫--拉盖尔多项式
\[
L_n(x) = \mathrm{e}^x \frac{\mathrm{d}^n}{\mathrm{d}x^n}(x^n \mathrm{e}^{-x})
\]
所有的根都是正数.
\end{prob}

\begin{prob}[M]{1243}
证明：切比雪夫--埃尔米特多项式
\[
H_n(x) = (-1)^n \mathrm{e}^{x^2} \frac{\mathrm{d}^n}{\mathrm{d}x^n}(\mathrm{e}^{-x^2})
\]
所有的根都是实数.
\end{prob}

\begin{prob}[S]{1244}
在曲线 $y = x^3$ 上某点的切线，平行于连接点 $A(-1, -1)$ 及点 $B(2, 8)$ 所成的弦，求出此点.
\end{prob}

\begin{prob}[S]{1245}
若 $ab < 0$，有限增量公式对于函数
\[
f(x) = \frac{1}{x}
\]
在闭区间 $[a, b]$ 上是否正确？
\end{prob}

\begin{prob}[L]{1246.1}
设
\begin{enumerate}[label=(\alph*), itemsep=2pt, topsep=3pt]
  \item $f(x) = ax^2 + bx + c$ $(a \neq 0)$;
  \item $f(x) = x^3$;
  \item $f(x) = \dfrac{1}{x}$;
  \item $f(x) = \mathrm{e}^x$,
\end{enumerate}
求满足
\[
f(x + \Delta x) - f(x) = \Delta x f'(x + \theta \Delta x) \quad (0 < \theta < 1)
\]
的函数 $\theta = \theta(x, \Delta x)$.
\end{prob}

\begin{prob}[M]{1246.2}
设 $f(x) \in C^{(1)}(-\infty, +\infty)$，且对于任何 $x$ 和 $h$ 成立恒等式
\[
f(x + h) - f(x) \equiv h f'(x).
\]
证明：
\[
f(x) = ax + b,
\]
其中 $a$ 和 $b$ 为常数.
\end{prob}

\begin{prob}[M]{1246.3}
设 $f(x) \in C^{(2)}(-\infty, +\infty)$，且对于任何 $x$ 和 $h$ 成立恒等式
\[
f(x + h) - f(x) \equiv h f'\left(x + \frac{h}{2}\right).
\]
证明：
\[
f(x) = ax^2 + bx + c,
\]
其中 $a, b$ 和 $c$ 为常数.
\end{prob}

\begin{prob}[M]{1247}
证明：若 $x \geqslant 0$，则
\[
\sqrt{x + 1} - \sqrt{x} = \frac{1}{2\sqrt{x + \theta(x)}},
\]
其中 $\dfrac{1}{4} \leqslant \theta(x) \leqslant \dfrac{1}{2}$，并且 $\lim\limits_{x \to 0} \theta(x) = \dfrac{1}{4}$，$\lim\limits_{x \to +\infty} \theta(x) = \dfrac{1}{2}$.
\end{prob}

\begin{prob}[M]{1248}
设
\[
f(x) = \begin{cases}
\dfrac{3 - x^2}{2}, & 0 \leqslant x \leqslant 1, \\[4pt]
\dfrac{1}{x}, & 1 < x < +\infty,
\end{cases}
\]
在闭区间 $[0, 2]$ 上对于函数 $f(x)$ 求有限增量公式中的中间值 $c$.
\end{prob}

\begin{prob}[M]{1249}
设
\[
f(x) - f(0) = x f'[\xi(x)],
\]
其中 $0 < \xi(x) < x$. 证明：若
\[
f(x) = \begin{cases}
x \sin(\ln x), & x > 0, \\
0, & x = 0,
\end{cases}
\]
则函数 $\xi = \xi(x)$ 在任意小的区间 $(0, \varepsilon)$ 内 $(\varepsilon > 0)$ 是不连续的.
\end{prob}

\begin{prob}[M]{1250}
设函数 $f(x)$ 在区间 $(a, b)$ 内有连续的导数 $f'(x)$. 对于区间 $(a, b)$ 内任何一点 $\xi$，可否从此区间中指出另外的两点 $x_1$ 及 $x_2$，使得
\[
\frac{f(x_2) - f(x_1)}{x_2 - x_1} = f'(\xi) \quad (x_1 < \xi < x_2)?
\]
研究一个例子：
\[
f(x) = x^3 \quad (-1 \leqslant x \leqslant 1), \quad \xi = 0.
\]
\end{prob}

\begin{prob}[L]{1251}
证明下列不等式：
\begin{enumerate}[label=(\alph*), itemsep=2pt, topsep=3pt]
  \item $|\sin x - \sin y| \leqslant |x - y|$;
  \item $p y^{p-1}(x - y) \leqslant x^p - y^p \leqslant p x^{p-1}(x - y) \quad (0 < y < x, p > 1)$;
  \item $|\arctan a - \arctan b| \leqslant |a - b|$;
  \item $\dfrac{a - b}{a} < \ln \dfrac{a}{b} < \dfrac{a - b}{b} \quad (0 < b < a)$.
\end{enumerate}
\end{prob}

\begin{prob}[S]{1252}
说明在闭区间 $[-1, 1]$ 上柯西定理对于函数 $f(x) = x^2$ 及 $g(x) = x^3$ 何以不真？
\end{prob}

\begin{prob}[M]{1253}
设函数 $f(x)$ 在闭区间 $[x_1, x_2]$ 上可微，并且 $x_1 x_2 > 0$. 证明：
\[
\frac{1}{x_1 - x_2}
\begin{vmatrix}
x_1 & x_2 \\
f(x_1) & f(x_2)
\end{vmatrix}
= f(\xi) - \xi f'(\xi),
\]
其中 $x_1 < \xi < x_2$.
\end{prob}

\begin{prob}[M]{1254}
证明：若函数 $f(x)$ 在有限的区间 $(a, b)$ 内可微但无界，则其导数 $f'(x)$ 在区间 $(a, b)$ 内也无界. 逆命题不真（举出例子）.
\end{prob}

\begin{prob}[M]{1255}
证明：若函数 $f(x)$ 在有限或无穷的区间 $(a, b)$ 内有有界的导数 $f'(x)$，则 $f(x)$ 在 $(a, b)$ 中一致连续.
\end{prob}

\begin{prob}[M]{1256}
证明：若函数 $f(x)$ 在无穷的区间 $(x_0, +\infty)$ 内可微，且
\[
\lim_{x \to +\infty} f'(x) = 0,
\]
则
\[
\lim_{x \to +\infty} \frac{f(x)}{x} = 0,
\]
即当 $x \to +\infty$ 时，$f(x) = o(x)$.
\end{prob}

\begin{prob}[M]{1257}
证明：若函数 $f(x)$ 在无穷的区间 $(x_0, +\infty)$ 内可微且
\[
\text{当 } x \to +\infty \text{ 时，} f(x) = o(x),
\]
则
\[
\lim_{x \to +\infty} |f'(x)| = 0.
\]
\end{prob}

\begin{prob}[M]{1258.1}
证明：若函数 $f(x)$：(1) 在闭区间 $[x_0, X]$ 上有定义并且是连续的；(2) 在区间 $(x_0, X)$ 内有有限的导数 $f'(x)$；(3) 存在有限或无穷的极限
\[
\lim_{x \to x_0 + 0} f'(x) = f'(x_0 + 0),
\]
则相应地存在有限或无穷的单侧导数 $f'_+(x_0)$ 且
\[
f'_+(x_0) = f'(x_0 + 0).
\]
\end{prob}

\begin{prob}[M]{1258.2}
证明：对于函数
\[
f(x) = \begin{cases}
\arctan \dfrac{1 + x}{1 - x}, & x \neq 1, \\
0, & x = 1
\end{cases}
\]
存在有限的极限
\[
\lim_{x \to 1} f'(x),
\]
但是函数 $f(x)$ 没有单侧导数 $f'_-(1)$ 及 $f'_+(1)$. 给出这个事实的几何解释. 然而，在这个点存在广义单侧导数（见习题 1009.2）.
\end{prob}

\begin{prob}[S]{1259}
证明：若当 $a < x < b$ 时 $f'(x) = 0$，则
\[
\text{当 } a < x < b \text{ 时，} f(x) = \text{常数}.
\]
\end{prob}

\begin{prob}[S]{1260}
证明：导数为常数
\[
f'(x) = k
\]
的唯一函数 $f(x)$ $(-\infty < x < +\infty)$ 是线性函数
\[
f(x) = kx + b.
\]
\end{prob}

\begin{prob}[S]{1261.1}
若 $f^{(n)}(x) = 0$，则函数 $f(x)$ 有什么性质？
\end{prob}

\begin{prob}[M]{1261.2}
设 $f(x) \in C^{(\infty)}(-\infty, +\infty)$，并且对每个 $x$ 都存在正整数 $n_x$ $(n_x \leqslant n)$，使得
\[
f^{(n_x)}(x) = 0.
\]
证明：函数 $f(x)$ 是多项式.
\end{prob}

\begin{prob}[M]{1262}
证明：满足方程
\[
y' = \lambda y \quad (\lambda = \text{常数})
\]
的唯一函数 $y = y(x)$ ($-\infty < x < +\infty$) 是指数函数 $y = C\mathrm{e}^{\lambda x}$，其中 $C$ 为任意常数.
\probhint{研究 $(y\mathrm{e}^{-\lambda x})'$.}
\end{prob}

\begin{prob}[M]{1263}
检验函数
\[
f(x) = \arctan \frac{1 + x}{1 - x}, \quad g(x) = \arctan x
\]
在区间：(1) $x < 1$ 及 (2) $x > 1$ 内有相同的导数. 推出这些函数间的关系.
\end{prob}

\begin{prob}[M]{1264}
证明下列恒等式：
\begin{enumerate}[label=(\alph*), itemsep=2pt, topsep=3pt]
\item $2\arctan x + \arcsin \dfrac{2x}{1 + x^2} = \pi \operatorname{sgn} x$，当 $|x| \geqslant 1$；
\item $3\arccos x - \arccos(3x - 4x^3) = \pi$，当 $|x| \leqslant \dfrac{1}{2}$.
\end{enumerate}
\end{prob}

\begin{prob}[M]{1265}
证明：若函数 $f(x)$：(1) 在闭区间 $[a, b]$ 上是连续的；(2) 于此区间内有有限的导数 $f'(x)$；(3) 不是线性函数，则在区间 $(a, b)$ 内至少能找到一点 $c$，使得
\[
|f'(c)| > \left|\frac{f(b) - f(a)}{b - a}\right|.
\]
给出这个事实的几何解释.
\end{prob}

\begin{prob}[M]{1266}
证明：若函数 $f(x)$：(1) 在闭区间 $[a, b]$ 上有二阶导数 $f''(x)$；(2) $f'(a) = f'(b) = 0$，则在区间 $(a, b)$ 内至少存在一点 $c$，使得
\[
|f''(c)| \geqslant \frac{4}{(b - a)^2}|f(b) - f(a)|.
\]
\end{prob}

\begin{prob}[M]{1267}
汽车从某点开始行驶，于 $t$ 秒内走完了路程，所经过的距离为 $s$. 证明：汽车运动的加速度在某瞬时不小于 $\dfrac{4s}{t^2}$.
\end{prob}

\probsec{§7. 增函数与减函数. 不等式}

\probgroup{求下列函数的严格单调 (增或减) 区间：}

\begin{prob}[S]{1268}
$y = 2 + x - x^2$.
\end{prob}

\begin{prob}[S]{1269}
$y = 3x - x^3$.
\end{prob}

\begin{prob}[S]{1270}
$y = \dfrac{2x}{1 + x^2}$.
\end{prob}

\begin{prob}[S]{1271}
$y = \dfrac{\sqrt{x}}{x + 100} \quad (x \geqslant 0)$.
\end{prob}

\begin{prob}[S]{1272}
$y = x + \sin x$.
\end{prob}

\begin{prob}[S]{1273}
$y = x + |\sin 2x|$.
\end{prob}

\begin{prob}[S]{1274}
$y = \cos \dfrac{\pi}{x}$.
\end{prob}

\begin{prob}[S]{1275}
$y = \dfrac{x^2}{2^x}$.
\end{prob}

\begin{prob}[S]{1276}
$y = x^n \mathrm{e}^{-x} \quad (n > 0, x \geqslant 0)$.
\end{prob}

\begin{prob}[S]{1277}
$y = x^2 - \ln x^2$.
\end{prob}

\begin{prob}[M]{1278}
若
\[
f(x) = \begin{cases}
x\left(\sqrt{\dfrac{3}{2}} + \sin\ln x\right), & x > 0, \\
0, & x = 0.
\end{cases}
\]
\end{prob}

\begin{prob}[M]{1279}
证明：圆的内接正 $n$ 边形的周长 $p_n$ 当边的数目 $n$ 增加时增加，而此圆的外切正 $n$ 边形的周长 $P_n$ 此时则减小. 利用这点来证明，当 $n \to \infty$ 时，$p_n$ 及 $P_n$ 有相同的极限.
\end{prob}

\begin{prob}[M]{1280}
证明：函数
\[
\left(1 + \frac{1}{x}\right)^x
\]
在区间 $(-\infty, -1)$ 及 $(0, +\infty)$ 内递增.
\end{prob}

\begin{prob}[M]{1281}
证明：有理函数
\[
P(x) = a_0 + a_1 x + \cdots + a_n x^n \quad (n \geqslant 1, a_n \neq 0)
\]
是区间 $(-\infty, -x_0)$ 及 $(x_0, +\infty)$ 上的严格单调函数，其中 $x_0$ 为充分大的正数.
\end{prob}

\begin{prob}[M]{1282}
证明：有理函数
\[
R(x) = \frac{a_0 + a_1 x + \cdots + a_n x^n}{b_0 + b_1 x + \cdots + b_m x^m} \quad (a_n b_m \neq 0)
\]
是区间 $(-\infty, -x_0)$ 及 $(x_0, +\infty)$ 上的严格单调函数，其中 $x_0$ 为充分大的正数.
\end{prob}

\begin{prob}[M]{1283}
单调函数的导数是否也必为单调的？研究一个例子：
\[
f(x) = x + \sin x.
\]
\end{prob}

\begin{prob}[M]{1284}
证明：若 $\varphi(x)$ 为可微的单调增函数，且当 $x \geqslant x_0$ 时，$|f'(x)| \leqslant \varphi'(x)$，则当 $x \geqslant x_0$ 时，$|f(x) - f(x_0)| \leqslant \varphi(x) - \varphi(x_0)$.
给出这个事实的几何解释.
\end{prob}

\begin{prob}[M]{1285}
设函数 $f(x)$ 在区间 $a \leqslant x < +\infty$ 内连续，而且当 $x > a$ 时，$f'(x) > k > 0$，其中 $k$ 为常数. 证明：若 $f(a) < 0$，则在区间 $\left(a, a - \dfrac{f(a)}{k}\right)$ 内方程 $f(x) = 0$ 有且仅有一个实根.
\end{prob}

\begin{prob}[M]{1286}
若在某邻域 $|x - x_0| < \delta$ 内，函数增量 $\Delta f(x_0) = f(x) - f(x_0)$ 的符号与自变量增量 $\Delta x_0 = x - x_0$ 的符号相同，则称函数 $f(x)$ 为在 $x_0$ 点的增函数.

证明：若函数 $f(x)$ ($a < x < b$) 在有限或无穷的区间 $(a, b)$ 内的每一点皆为增函数，则它在此区间内为增函数.
\end{prob}

\begin{prob}[M]{1287}
证明：函数
\[
f(x) = \begin{cases}
x + x^2 \sin\dfrac{2}{x}, & x \neq 0, \\
0, & x = 0
\end{cases}
\]
在点 $x = 0$ 是增函数，但在包含此点的任何区间 $(-\varepsilon, \varepsilon)$ 中并非增函数，其中 $\varepsilon > 0$ 为任意小的数. 作出此函数的略图.
\end{prob}

\begin{prob}[M]{1288}
证明定理：设 (1) 函数 $\varphi(x)$ 及 $\psi(x)$ 为 $n$ 阶可微函数；(2) $\varphi^{(k)}(x_0) = \psi^{(k)}(x_0)$ ($k = 0, 1, \cdots, n - 1$)；(3) 当 $x > x_0$ 时，$\varphi^{(n)}(x) > \psi^{(n)}(x)$，则当 $x > x_0$ 时有不等式
\[
\varphi(x) > \psi(x).
\]
\end{prob}

\begin{prob}[L]{1289}
证明下列不等式：
\begin{enumerate}[label=(\alph*), itemsep=2pt, topsep=3pt]
\item $\mathrm{e}^x > 1 + x \quad (x \neq 0)$；
\item $x - \dfrac{x^2}{2} < \ln(1 + x) < x \quad (x > 0)$；
\item $x - \dfrac{x^3}{6} < \sin x < x \quad (x > 0)$；
\item $\tan x > x + \dfrac{x^3}{3} \quad \left(0 < x < \dfrac{\pi}{2}\right)$；
\item $(x^\alpha + y^\alpha)^{\frac{1}{\alpha}} > (x^\beta + y^\beta)^{\frac{1}{\beta}} \quad (x > 0, y > 0, 0 < \alpha < \beta)$.
\end{enumerate}
给出不等式 (a)---(d) 的几何解释.
\end{prob}

\begin{prob}[M]{1290}
证明：不等式
\[
\frac{2}{\pi} x < \sin x < x
\]
当 $0 < x < \dfrac{\pi}{2}$ 时成立.
\end{prob}

\begin{prob}[M]{1291}
证明：当 $x > 0$ 时有不等式
\[
\left(1 + \frac{1}{x}\right)^x < \mathrm{e} < \left(1 + \frac{1}{x}\right)^{x+1}.
\]
\end{prob}

\begin{prob}[M]{1292}
等差数列与等比数列的项的数目相同且有相同的首项与末项，并且一切项都是正的. 证明：等差数列各项之和大于等于等比数列各项之和.
\end{prob}

\begin{prob}[M]{1293}
用不等式
\[
\sum_{k=1}^{n} (a_k x + b_k)^2 \geqslant 0,
\]
其中 $x, a_k, b_k$ ($k = 1, \cdots, n$) 为实数，来证明柯西--布尼亚科夫斯基不等式：
\[
\left(\sum_{k=1}^{n} a_k b_k\right)^2 \leqslant \sum_{k=1}^{n} a_k^2 \cdot \sum_{k=1}^{n} b_k^2.
\]
\end{prob}

\begin{prob}[M]{1294}
证明：正数的算术平均值的平方不大于这些数的平方平均值，即
\[
\frac{1}{n}\sum_{k=1}^{n} x_k \leqslant \sqrt{\frac{1}{n}\sum_{k=1}^{n} x_k^2}.
\]
\end{prob}

\begin{prob}[M]{1295}
证明：正数的几何平均值不大于这些数的算术平均值，即
\[
\sqrt[n]{x_1 x_2 \cdots x_n} \leqslant \frac{1}{n}(x_1 + x_2 + \cdots + x_n).
\]
\probhint{用数学归纳法.}
\end{prob}

\begin{prob}[L]{1296}
设 $a$ 及 $b$ 为两个正数，则由等式
\[
\Delta_s(a, b) = \begin{cases}
\left(\dfrac{a^s + b^s}{2}\right)^{\frac{1}{s}}, & s \neq 0, \\
\lim\limits_{t \to 0} \Delta_t(a, b), & s = 0
\end{cases}
\]
所定义的函数称为正数 $a$ 与 $b$ 之 $s$ 阶平均值.

例如，当 $s = -1$ 时得调和平均值，当 $s = 0$ 时得几何平均值 (试证明！)；当 $s = 1$ 时得算术平均值；当 $s = 2$ 时得平方平均值.

证明：(1) $\min(a, b) \leqslant \Delta_s(a, b) \leqslant \max(a, b)$；
(2) 当 $a \neq b$ 时，函数 $\Delta_s(a, b)$ 是变量 $s$ 的增函数；
(3) $\lim\limits_{s \to -\infty} \Delta_s(a, b) = \min(a, b)$，$\lim\limits_{s \to +\infty} \Delta_s(a, b) = \max(a, b)$.
\probhint{考察 $\dfrac{\mathrm{d}}{\mathrm{d}s}[\ln \Delta_s(a, b)]$.}
\end{prob}

\begin{prob}[L]{1297}
证明下列不等式：
\begin{enumerate}[label=(\alph*), itemsep=2pt, topsep=3pt]
\item $x^\alpha - 1 > \alpha(x - 1)$，若 $\alpha \geqslant 2, x > 1$；
\item $\sqrt[n]{x} - \sqrt[n]{a} < \sqrt[n]{x - a}$，若 $n > 1, x > a > 0$；
\item $1 + 2\ln x \leqslant x^2$，若 $x > 0$.
\end{enumerate}
\end{prob}

\probsec{§8. 凹凸性. 拐点}

\begin{prob}[M]{1298}
研究曲线
\[
y = 1 + \sqrt[3]{x}
\]
于 $A(-1, 0)$，$B(1, 2)$ 及 $C(0, 1)$ 诸点的凹凸性.
\end{prob}

\probgroup{求下列函数的图像的凹或凸的区间及拐点：}

\begin{prob}[S]{1299}
$y = 3x^2 - x^3$.
\end{prob}

\begin{prob}[S]{1300}
$y = \dfrac{a^2}{a^2 + x^2} \quad (a > 0)$.
\end{prob}

\begin{prob}[S]{1301}
$y = x + x^{\frac{5}{3}}$.
\end{prob}

\begin{prob}[S]{1302}
$y = \sqrt{1 + x^2}$.
\end{prob}

\begin{prob}[S]{1303}
$y = x + \sin x$.
\end{prob}

\begin{prob}[S]{1304}
$y = \mathrm{e}^{-x^2}$.
\end{prob}

\begin{prob}[S]{1305}
$y = \ln(1 + x^2)$.
\end{prob}

\begin{prob}[S]{1306}
$y = x \sin(\ln x) \quad (x > 0)$.
\end{prob}

\begin{prob}[S]{1307}
$y = x^x \quad (x > 0)$.
\end{prob}

\begin{prob}[M]{1308}
证明：曲线
\[
y = \frac{x + 1}{x^2 + 1}
\]
在位于同一直线上的三个拐点. 作出这个函数的图像.
\end{prob}

\begin{prob}[M]{1309}
如何选择参量 $h$，可使“概率曲线”
\[
y = \frac{h}{\sqrt{\pi}} \mathrm{e}^{-h^2 x^2} \quad (h > 0)
\]
有拐点 $x = \pm \sigma$？
\end{prob}

\begin{prob}[M]{1310}
研究摆线 (旋轮线)
\[
x = a(t - \sin t), \quad y = a(1 - \cos t) \quad (a > 0)
\]
的凹凸性.
\end{prob}

\begin{prob}[M]{1311}
设函数 $f(x)$ 在区间 $a \leqslant x < +\infty$ 中二阶可微，并且：
(1) $f(a) = A > 0$； (2) $f'(a) < 0$； (3) 当 $x > a$ 时，$f''(x) \leqslant 0$.
证明：在区间 $(a, +\infty)$ 内方程 $f(x) = 0$ 有且仅有一个实根.
\end{prob}

\begin{prob}[M]{1312}
若对于区间 $(a, b)$ 内的任意两点 $x_1$ 与 $x_2$ 及任意两数 $\lambda_1$ 与 $\lambda_2$ ($\lambda_1 > 0, \lambda_2 > 0, \lambda_1 + \lambda_2 = 1$) 有不等式：
\[
f(\lambda_1 x_1 + \lambda_2 x_2) < \lambda_1 f(x_1) + \lambda_2 f(x_2)
\]
(或对应地，相反的不等式 $f(\lambda_1 x_1 + \lambda_2 x_2) > \lambda_1 f(x_1) + \lambda_2 f(x_2)$)，则称函数 $f(x)$ 在区间 $(a, b)$ 上是凸 (凹) 的.

证明：(1) 若当 $a < x < b$ 时，$f''(x) > 0$，则函数 $f(x)$ 在 $(a, b)$ 上是凸的；(2) 若当 $a < x < b$ 时，$f''(x) < 0$，则 $f(x)$ 在 $(a, b)$ 上是凹的.
\end{prob}

\begin{prob}[M]{1313}
证明：函数
\[
x^n \ (n > 1), \quad \mathrm{e}^x, \quad x \ln x
\]
在区间 $(0, +\infty)$ 上是凸的，而函数
\[
x^n \ (0 < n < 1), \quad \ln x
\]
在区间 $(0, +\infty)$ 上是凹的.
\end{prob}

\begin{prob}[L]{1314.1}
证明下列不等式，并解释其几何意义：
\begin{enumerate}[label=(\alph*), itemsep=2pt, topsep=3pt]
\item $\dfrac{1}{2}(x^n + y^n) > \left(\dfrac{x + y}{2}\right)^n \ (x > 0, y > 0, x \neq y, n > 1)$；
\item $\dfrac{\mathrm{e}^x + \mathrm{e}^y}{2} > \mathrm{e}^{\frac{x+y}{2}} \ (x \neq y)$；
\item $x \ln x + y \ln y > (x + y)\ln\dfrac{x + y}{2} \ (x > 0, y > 0)$.
\end{enumerate}
\end{prob}

\begin{prob}[M]{1314.2}
设当 $a \leqslant x \leqslant b$ 时 $f''(x) \geqslant 0$，证明：对于任何 $x_1, x_2 \in [a, b]$，都有
\[
f\left(\frac{x_1 + x_2}{2}\right) \leqslant \frac{1}{2}[f(x_1) + f(x_2)].
\]
\end{prob}

\begin{prob}[M]{1315}
证明：有界的凸函数处处连续，并有左导数和右导数.
\end{prob}

\begin{prob}[M]{1316}
设函数 $f(x)$ 在区间 $(a, b)$ 内二阶可微，且 $f''(\xi) \neq 0$，其中 $a < \xi < b$. 证明：在区间 $(a, b)$ 中可找出两个值 $x_1$ 与 $x_2$，满足
\[
\frac{f(x_2) - f(x_1)}{x_2 - x_1} = f'(\xi).
\]
\end{prob}

\begin{prob}[M]{1317}
证明：若函数 $f(x)$ 在无穷的区间 $(x_0, +\infty)$ 内二阶可微，且
\[
\lim_{x \to x_0 + 0} f(x) = 0, \quad \lim_{x \to +\infty} f(x) = 0,
\]
则在区间 $(x_0, +\infty)$ 内最少有一点 $\xi$，满足 $f''(\xi) = 0$.
\end{prob}

\probsec{§9. 不定式的求值法}

\probgroup{求出下列各式之值：}

\begin{prob}[S]{1318}
$\lim_{x \to 0} \frac{\sin ax}{\sin bx}$.
\end{prob}

\begin{prob}[S]{1319}
$\lim_{x \to 0} \frac{\cosh x - \cos x}{x^2}$.
\end{prob}

\begin{prob}[S]{1320}
$\lim_{x \to 0} \frac{\tan x - x}{x - \sin x}$.
\end{prob}

\begin{prob}[S]{1321}
$\lim_{x \to 0} \frac{3\tan 4x - 12\tan x}{3\sin 4x - 12\sin x}$.
\end{prob}

\begin{prob}[S]{1322}
$\lim_{x \to \frac{\pi}{2}} \frac{\tan 3x}{\tan x}$.
\end{prob}

\begin{prob}[S]{1323}
$\lim_{x \to 0} \frac{x\cot x - 1}{x^2}$.
\end{prob}

\begin{prob}[S]{1324}
$\lim_{x \to \frac{\pi}{4}} \frac{\sqrt[3]{\tan x} - 1}{2\sin^2 x - 1}$.
\end{prob}

\begin{prob}[S]{1325}
$\lim_{x \to 0} \frac{x(\mathrm{e}^x + 1) - 2(\mathrm{e}^x - 1)}{x^3}$.
\end{prob}

\begin{prob}[S]{1326}
$\lim_{x \to 0} \frac{1 - \cos x^2}{x^2 \sin x^2}$.
\end{prob}

\begin{prob}[S]{1327}
$\lim_{x \to 0} \frac{\arcsin 2x - 2\arcsin x}{x^3}$.
\end{prob}

\begin{prob}[S]{1328}
$\lim_{x \to 0} \frac{1}{x\sqrt{x}} \left( \sqrt{a}\arctan\sqrt{\frac{x}{a}} - \sqrt{b}\arctan\sqrt{\frac{x}{b}} \right)$.
\end{prob}

\begin{prob}[S]{1329}
$\lim_{x \to 0} \frac{a^x - a^{\sin x}}{x^3} \quad (a > 0)$.
\end{prob}

\begin{prob}[S]{1330}
$\lim_{x \to 1} \left( \frac{x^x - x}{\ln x - x + 1} \right)$.
\end{prob}

\begin{prob}[S]{1331}
$\lim_{x \to 0} \frac{\ln(\sin ax)}{\ln(\sin bx)}$.
\end{prob}

\begin{prob}[S]{1332}
$\lim_{x \to 0} \frac{\ln(\cos ax)}{\ln(\cos bx)}$.
\end{prob}

\begin{prob}[S]{1333}
$\lim_{x \to 0} \frac{\cos(\sin x) - \cos x}{x^4}$.
\end{prob}

\begin{prob}[S]{1334}
$\lim_{x \to 0} \frac{1}{x} \left( \frac{1}{\tanh x} - \frac{1}{\tan x} \right)$.
\end{prob}

\begin{prob}[S]{1335}
$\lim_{x \to 0} \frac{\operatorname{arsinh}(\sinh x) - \operatorname{arsinh}(\sin x)}{\sinh x - \sin x}$，其中 $\operatorname{arsinh} x = \ln(x + \sqrt{1 + x^2})$。
\end{prob}

\begin{prob}[S]{1336}
$\lim_{x \to +\infty} \frac{\ln x}{x^\varepsilon} \quad (\varepsilon > 0)$.
\end{prob}

\begin{prob}[S]{1337}
$\lim_{x \to +\infty} \frac{x^n}{\mathrm{e}^{ax}} \quad (a > 0, n > 0)$.
\end{prob}

\begin{prob}[S]{1338}
$\lim_{x \to 0} \frac{\mathrm{e}^{-\frac{1}{x^2}}}{x^{1000}}$.
\end{prob}

\begin{prob}[S]{1339}
$\lim_{x \to +\infty} x^2 \mathrm{e}^{-0.01x}$.
\end{prob}

\begin{prob}[S]{1340}
$\lim_{x \to 1-0} \ln x \cdot \ln(1 - x)$.
\end{prob}

\begin{prob}[S]{1341}
$\lim_{x \to +0} x^\varepsilon \ln x \quad (\varepsilon > 0)$.
\end{prob}

\begin{prob}[S]{1342}
$\lim_{x \to +0} x^x$.
\end{prob}

\begin{prob}[S]{1343}
$\lim_{x \to 0} x^{x^x - 1}$.
\end{prob}

\begin{prob}[S]{1344}
$\lim_{x \to 0} \left( x^{x^x} - 1 \right)$.
\end{prob}

\begin{prob}[S]{1345}
$\lim_{x \to +0} x^{\frac{k}{1 + \ln x}}$.
\end{prob}

\begin{prob}[S]{1346}
$\lim_{x \to 1} x^{\frac{1}{1 - x}}$.
\end{prob}

\begin{prob}[S]{1347}
$\lim_{x \to 1} (2 - x)^{\tan\frac{\pi x}{2}}$.
\end{prob}

\begin{prob}[S]{1348}
$\lim_{x \to \frac{\pi}{4}} (\tan x)^{\tan 2x}$.
\end{prob}

\begin{prob}[S]{1349}
$\lim_{x \to 0} (\cot x)^{\sin x}$.
\end{prob}

\begin{prob}[S]{1350}
$\lim_{x \to +0} \left( \ln\frac{1}{x} \right)^x$.
\end{prob}

\begin{prob}[S]{1351}
$\lim_{x \to \infty} \left( \tan\frac{\pi x}{2x + 1} \right)^{\frac{1}{x}}$.
\end{prob}

\begin{prob}[S]{1352}
$\lim_{x \to a} \left( \frac{\tan x}{\tan a} \right)^{\cot(x - a)}$.
\end{prob}

\begin{prob}[S]{1353}
$\lim_{x \to 0} \left( \frac{a^x - x\ln a}{b^x - x\ln b} \right)^{\frac{1}{x^2}}$.
\end{prob}

\begin{prob}[S]{1354}
$\lim_{x \to 0} \left( \frac{1}{x} - \frac{1}{\mathrm{e}^x - 1} \right)$.
\end{prob}

\begin{prob}[S]{1355}
$\lim_{x \to 1} \left( \frac{1}{\ln x} - \frac{1}{x - 1} \right)$.
\end{prob}

\begin{prob}[S]{1356}
$\lim_{x \to 0} \left( \cot x - \frac{1}{x} \right)$.
\end{prob}

\begin{prob}[S]{1357}
$\lim_{x \to 0} \left[ \frac{1}{\ln(x + \sqrt{1 + x^2})} - \frac{1}{\ln(1 + x)} \right]$.
\end{prob}

\begin{prob}[S]{1358}
$\lim_{x \to a} \frac{a^x - x^a}{x - a} \quad (a > 0)$.
\end{prob}

\begin{prob}[S]{1359}
$\lim_{x \to 0} \frac{(1 + x)^{\frac{1}{x}} - \mathrm{e}}{x}$.
\end{prob}

\begin{prob}[S]{1360}
$\lim_{x \to 0} \frac{(a + x)^x - a^x}{x^2} \quad (a > 0)$.
\end{prob}

\begin{prob}[S]{1361}
$\lim_{x \to +\infty} \left( \frac{2}{\pi}\arctan x \right)^x$.
\end{prob}

\begin{prob}[S]{1362}
$\lim_{x \to +\infty} (\tanh x)^x$.
\end{prob}

\begin{prob}[L]{1363}
\begin{enumerate}[label=(\alph*), itemsep=2pt, topsep=3pt]
\item $\lim_{x \to 0} \left( \frac{\arcsin x}{x} \right)^{\frac{1}{x^2}}$;
\item $\lim_{x \to 0} \left( \frac{\sin x}{x} \right)^{\frac{1}{x^2}}$;
\item $\lim_{x \to 0} \left( \frac{\tan x}{x} \right)^{\frac{1}{x^2}}$;
\item $\lim_{x \to 0} \left( \frac{\arctan x}{x} \right)^{\frac{1}{x^2}}$;
\item $\lim_{x \to 0} \left( \frac{\operatorname{arsinh} x}{x} \right)^{\frac{1}{x^2}}$，其中 $\operatorname{arsinh} x = \ln(x + \sqrt{1 + x^2})$。
\end{enumerate}
\end{prob}

\begin{prob}[S]{1364}
$\lim_{x \to 0} \left[ \frac{(1 + x)^{\frac{1}{x}}}{\mathrm{e}} \right]^{\frac{1}{x}}$.
\end{prob}

\begin{prob}[S]{1365}
$\lim_{x \to 0} \left( \frac{2}{\pi}\arccos x \right)^{\frac{1}{x}}$.
\end{prob}

\begin{prob}[S]{1366}
$\lim_{x \to 0} \left( \frac{\cos x}{\cosh x} \right)^{\frac{1}{x^2}}$.
\end{prob}

\begin{prob}[S]{1367}
$\lim_{x \to 0} \frac{\ln\cosh x}{\sqrt[m]{\cosh x} - \sqrt[n]{\cosh x}}$.
\end{prob}

\begin{prob}[M]{1368}
(a) $\lim_{x \to 0} \left( \frac{1 + \mathrm{e}^x}{2} \right)^{\coth x}$; \qquad (b) $\lim_{x \to +\infty} \frac{x^{\ln x}}{(\ln x)^x}$.
\end{prob}

\begin{prob}[S]{1369}
$\lim_{x \to +\infty} \left[ \sqrt[3]{x^3 + x^2 + x + 1} - \sqrt{x^2 + x + 1} \cdot \frac{\ln(\mathrm{e}^x + x)}{x} \right]$.
\end{prob}

\begin{prob}[S]{1370}
$\lim_{x \to +\infty} \left[ (x + a)^{1 + \frac{1}{x}} - x^{1 + \frac{1}{x + a}} \right]$.
\end{prob}

\begin{prob}[M]{1371}
若当 $x \to 0$ 时，曲线 $y = f(x)$ 通过坐标原点 $(0, 0)$ [$\lim_{x \to 0} f(x) = 0$]，且在此有倾角 $\alpha$，求 $\lim_{x \to 0} \frac{y}{x}$。
\end{prob}

\begin{prob}[M]{1372}
若当 $x \to +0$ 时，连续曲线 $y = f(x)$ 通过坐标原点 $(0, 0)$ [$\lim_{x \to +0} f(x) = 0$]，并且当 $0 < x < \varepsilon$ 时，此曲线完全位于两直线 $y = -kx$ 及 $y = kx$ ($k \neq \infty$) 所组成的锐角之内，证明：$\lim_{x \to +0} x^{f(x)} = 1$。
\end{prob}

\begin{prob}[M]{1373.1}
证明：若函数 $f(x)$ 的二阶导数 $f''(x)$ 存在，则
\[
f''(x) = \lim_{h \to 0} \frac{f(x + h) + f(x - h) - 2f(x)}{h^2}.
\]
\end{prob}

\begin{prob}[M]{1373.2}
研究函数
\[
f(x) = \begin{cases} \dfrac{1}{x} - \dfrac{1}{\mathrm{e}^x - 1}, & x \neq 0, \\ \dfrac{1}{2}, & x = 0 \end{cases}
\]
在点 $x = 0$ 的可微性。
\end{prob}

\begin{prob}[M]{1373.3}
求曲线
\[
y = \frac{x^{1 + x}}{(1 + x)^x} \quad (x > 0)
\]
的渐近线。
\end{prob}

\begin{prob}[L]{1374}
研究运用洛必达法则于下列各例的可能性：
\begin{enumerate}[label=(\alph*), itemsep=2pt, topsep=3pt]
\item $\lim_{x \to 0} \frac{x^2 \sin\frac{1}{x}}{\sin x}$;
\item $\lim_{x \to \infty} \frac{x - \sin x}{x + \sin x}$;
\item $\lim_{x \to +\infty} \frac{\mathrm{e}^{-2x}(\cos x + 2\sin x) + \mathrm{e}^{-x^2}\sin^2 x}{\mathrm{e}^{-x}(\cos x + \sin x)}$;
\item $\lim_{x \to \infty} \frac{1 + x + \sin x\cos x}{(x + \sin x\cos x)\mathrm{e}^{\sin x}}$.
\end{enumerate}
\end{prob}

\begin{prob}[M]{1375}
设有一弓形，其弦长为 $b$，拱高为 $h$，半径为 $R$，又有内接于此弓形的等腰三角形。若当 $R$ 不变时弓形的弧长趋于零，求弓形面积与内接等腰三角形面积之比的极限。利用所得结果推出弓形面积的近似公式：
\[
S \approx \frac{2}{3} bh.
\]
\end{prob}

\probsec{§10. 泰勒公式}

\begin{prob}[M]{1376}
将多项式
\[
P(x) = 1 + 3x + 5x^2 - 2x^3
\]
改写为二项式 $x + 1$ 的非负整数次幂多项式。
\end{prob}

\probgroup{写出下列函数按变量 $x$ 的非负整数次幂的展开式，到含有所指阶数的项为止：}

\begin{prob}[S]{1377}
$\frac{1 + x + x^2}{1 - x + x^2}$ 到含 $x^4$ 的项。$f^{(4)}(0)$ 等于什么？
\end{prob}

\begin{prob}[S]{1378}
$\frac{(1 + x)^{100}}{(1 - 2x)^{40}(1 + 2x)^{60}}$ 到含 $x^2$ 的项。
\end{prob}

\begin{prob}[S]{1379}
$\sqrt[m]{a^m + x} \quad (a > 0)$ 到含 $x^2$ 的项。
\end{prob}

\begin{prob}[S]{1380}
$\sqrt{1 - 2x + x^3} - \sqrt[3]{1 - 3x + x^2}$ 到含 $x^3$ 的项。
\end{prob}

\begin{prob}[S]{1381}
$\mathrm{e}^{2x - x^2}$ 到含 $x^5$ 的项。
\end{prob}

\begin{prob}[S]{1382}
$\frac{x}{\mathrm{e}^x - 1}$ 到含 $x^4$ 的项。
\end{prob}

\begin{prob}[S]{1383}
$\sqrt[3]{\sin x^3}$ 到含 $x^{13}$ 的项。
\end{prob}

\begin{prob}[S]{1384}
$\ln\cos x$ 到含 $x^6$ 的项。
\end{prob}

\begin{prob}[S]{1385}
$\sin(\sin x)$ 到含 $x^3$ 的项。
\end{prob}

\begin{prob}[S]{1386}
$\tan x$ 到含 $x^5$ 的项。
\end{prob}

\begin{prob}[S]{1387}
$\ln\frac{\sin x}{x}$ 到含 $x^6$ 的项。
\end{prob}

\begin{prob}[M]{1388}
求函数 $f(x) = \sqrt{x}$ 按照 $x - 1$ 的非负整数次幂展开式的前三项。
\end{prob}

\begin{prob}[M]{1389}
将函数 $f(x) = x^x - 1$ 按照 $x - 1$ 的非负整数次幂展开至含有 $(x - 1)^3$ 的项。
\end{prob}

\begin{prob}[M]{1390}
在点 $x = 0$ 的邻域中，用抛物线（二次多项式）近似地代替函数
\[
y = a\cosh\frac{x}{a} \quad (a > 0).
\]
\end{prob}

\begin{prob}[M]{1391}
按公式 $\frac{1}{x}$ 的非负整数次幂展开函数 $f(x) = \sqrt{1 + x^2} - x$ $(x > 0)$ 到含 $\frac{1}{x^3}$ 的项。
\end{prob}

\begin{prob}[M]{1392}
求函数 $f(h) = \ln(x + h)$ $(x > 0)$ 按增量 $h$ 的正整数次幂的展开式，到含 $h^n$ 的项（$n$ 为正整数）。
\end{prob}

\begin{prob}[M]{1393.1}
设
\[
f(x + h) = f(x) + hf'(x) + \cdots + \frac{h^n}{n!} f^{(n)}(x + \theta h) \quad (0 < \theta < 1),
\]
且 $f^{(n+1)}(x) \neq 0$。证明：
\[
\lim_{h \to 0} \theta = \frac{1}{n + 1}.
\]
\end{prob}

\begin{prob}[M]{1393.2}
设当 $x \to 0$ 时有 $f(x) = 1 + kx + o(x)$。证明：
\[
\lim_{x \to 0} [f(x)]^{\frac{1}{x}} = \mathrm{e}^k.
\]
\end{prob}

\begin{prob}[M]{1393.3}
设 $f(x) \in C^{(2)}[0, 1]$，$f(0) = f(1) = 0$，且当 $x \in (0, 1)$ 时 $|f''(x)| \leqslant A$。证明：当 $0 \leqslant x \leqslant 1$ 时
\[
|f'(x)| \leqslant \frac{A}{2}.
\]
\end{prob}

\begin{prob}[M]{1393.4}
设 $f(x)$ $(-\infty < x < +\infty)$ 为二阶可微函数，且
\[
M_k = \sup_{-\infty < x < +\infty} |f^{(k)}(x)| < +\infty \quad (k = 0, 1, 2).
\]
证明不等式:
\[
M_1^2 \leqslant 2M_0 M_2.
\]
\end{prob}

\begin{prob}[L]{1394}
估计下列近似公式的绝对误差：
\begin{enumerate}[label=(\alph*), itemsep=2pt, topsep=3pt]
    \item $\mathrm{e}^x \approx 1 + x + \dfrac{x^2}{2!} + \cdots + \dfrac{x^n}{n!}$，当 $0 \leqslant x \leqslant 1$；
    \item $\sin x \approx x - \dfrac{x^3}{6}$，当 $|x| \leqslant \dfrac{1}{2}$；
    \item $\tan x \approx x + \dfrac{x^3}{3}$，当 $|x| \leqslant 0.1$；
    \item $\sqrt{1 + x} \approx 1 + \dfrac{x}{2} - \dfrac{x^2}{8}$，当 $0 \leqslant x \leqslant 1$.
\end{enumerate}
\end{prob}

\begin{prob}[M]{1395.1}
近似公式
\[
\cos x = 1 - \frac{x^2}{2}
\]
对于怎样的 $x$ 精确到 $0.0001$?
\end{prob}

\begin{prob}[M]{1395.2}
证明公式：
\[
\sqrt[n]{a^n + x} = a + \frac{x}{na^{n-1}} - r \quad (n \geqslant 2, \, a > 0, \, x > 0),
\]
其中 $0 < r < \dfrac{n - 1}{2n^2} \cdot \dfrac{x^2}{a^{2n - 1}}$.
\end{prob}

\begin{prob}[L]{1396}
利用泰勒公式近似地计算并估计误差：
\begin{enumerate}[label=(\alph*), itemsep=2pt, topsep=3pt]
    \item $\sqrt[3]{30}$；
    \item $\sqrt[5]{250}$；
    \item $\sqrt[12]{4000}$；
    \item $\sqrt{\mathrm{e}}$；
    \item $\sin 18^\circ$；
    \item $\ln 1.2$；
    \item $\arctan 0.8$；
    \item $\arcsin 0.45$；
    \item $1.1^{1.2}$.
\end{enumerate}
\end{prob}

\begin{prob}[L]{1397}
计算：
\begin{enumerate}[label=(\alph*), itemsep=2pt, topsep=3pt]
    \item $\mathrm{e}$，精确到 $10^{-9}$；
    \item $\sin 1^\circ$，精确到 $10^{-8}$；
    \item $\cos 9^\circ$，精确到 $10^{-5}$；
    \item $\sqrt{5}$，精确到 $10^{-4}$；
    \item $\lg 11$，精确到 $10^{-5}$.
\end{enumerate}
\end{prob}

\probgroup{利用展开式 $\mathrm{I}\text{---}\mathrm{V}$，求下列极限：}

\begin{prob}[S]{1398}
$\lim\limits_{x \to 0} \dfrac{\cos x - \mathrm{e}^{-\frac{x^2}{2}}}{x^4}$.
\end{prob}

\begin{prob}[S]{1399}
$\lim\limits_{x \to 0} \dfrac{\mathrm{e}^x \sin x - x(1 + x)}{x^3}$.
\end{prob}

\begin{prob}[S]{1400}
$\lim\limits_{x \to +\infty} x^{\frac{3}{2}} (\sqrt{x + 1} + \sqrt{x - 1} - 2\sqrt{x})$.
\end{prob}

\begin{prob}[S]{1401}
$\lim\limits_{x \to +\infty} (\sqrt[6]{x^6 + x^5} - \sqrt[6]{x^6 - x^5})$.
\end{prob}

\begin{prob}[S]{1402}
$\lim\limits_{x \to +\infty} \left[\left(x^3 - x^2 + \dfrac{x}{2}\right)\mathrm{e}^{\frac{1}{x}} - \sqrt{x^6 + 1}\right]$.
\end{prob}

\begin{prob}[S]{1403}
$\lim\limits_{x \to 0} \dfrac{a^x + a^{-x} - 2}{x^2} \ (a > 0)$.
\end{prob}

\begin{prob}[S]{1404}
$\lim\limits_{x \to \infty} \left[x - x^2 \ln\left(1 + \dfrac{1}{x}\right)\right]$.
\end{prob}

\begin{prob}[S]{1405}
$\lim\limits_{x \to 0} \left(\dfrac{1}{x} - \dfrac{1}{\sin x}\right)$.
\end{prob}

\begin{prob}[L]{1406}
\begin{enumerate}[label=(\alph*), itemsep=2pt, topsep=3pt]
    \item $\lim\limits_{x \to 0} \dfrac{1}{x} \left(\dfrac{1}{x} - \cot x\right)$；
    \item $\lim\limits_{x \to 0} \dfrac{\sin(\sin x) - x\sqrt[3]{1 - x^2}}{x^5}$；
    \item $\lim\limits_{x \to 0} \dfrac{1 - (\cos x)^{\sin x}}{x^3}$；
    \item $\lim\limits_{x \to 0} \dfrac{\sinh(\tan x) - x}{x^3}$.
\end{enumerate}
\end{prob}

\probgroup{当 $x \to 0$ 时，求出无穷小量 $y$ 的形如 $Cx^n$（$C$ 为常数）的主项，设：}

\begin{prob}[S]{1407}
$y = \tan(\sin x) - \sin(\tan x)$.
\end{prob}

\begin{prob}[S]{1408}
$y = (1 + x)^x - 1$.
\end{prob}

\begin{prob}[S]{1409}
$y = 1 - \dfrac{(1 + x)^{\frac{1}{x}}}{\mathrm{e}}$.
\end{prob}

\begin{prob}[M]{1410.1}
当选择怎样的系数 $a$ 与 $b$ 时，量
\[
x - (a + b\cos x)\sin x
\]
对于 $x$ 为 $5$ 阶无穷小?
\end{prob}

\begin{prob}[M]{1410.2}
选择系数 $A$ 和 $B$，使当 $x \to 0$ 时成立渐近公式
\[
\cot x = \frac{1 + Ax^2}{x + Bx^3} + O(x^5).
\]
\end{prob}

\begin{prob}[M]{1410.3}
系数 $A, B, C, D$ 为何值时，在 $x \to 0$ 时有渐近公式
\[
\mathrm{e}^x = \frac{1 + Ax + Bx^2}{1 + Cx + Dx^2} + O(x^5)\text{?}
\]
\end{prob}

\begin{prob}[L]{1411}
假设 $|x|$ 为小量，推出下列各式的简单的近似公式：
\begin{enumerate}[label=(\alph*), itemsep=2pt, topsep=3pt]
    \item $\dfrac{1}{R^2} - \dfrac{1}{(R + x)^2} \ (R > 0)$；
    \item $\sqrt[3]{\dfrac{1 + x}{1 - x}} - \sqrt[3]{\dfrac{1 - x}{1 + x}}$；
    \item $\dfrac{A}{x}\left[1 - \left(1 + \dfrac{x}{100}\right)^{-n}\right]$；
    \item $\dfrac{\ln 2}{\ln\left(1 + \dfrac{x}{100}\right)}$.
\end{enumerate}
\end{prob}

\begin{prob}[M]{1412}
假设 $x$ 的绝对值为小量，推出形如
\[
x = \alpha \sin x + \beta \tan x
\]
且精确到 $x^5$ 项的近似公式. 应用此公式近似地求小角度弧长.
\end{prob}

\begin{prob}[M]{1413}
估计下面的切比雪夫法则的相对误差：圆弧长近似地等于以此弧的弦为底、以其拱高的 $\sqrt{\dfrac{4}{3}}$ 为高的等腰三角形的两腰之和.
\end{prob}

\probsec{§11. 函数的极值. 函数的最大值和最小值}

\probgroup{研究下列函数的极值：}

\begin{prob}[S]{1414}
$y = 2 + x - x^2$.
\end{prob}

\begin{prob}[S]{1415}
$y = (x - 1)^3$.
\end{prob}

\begin{prob}[S]{1416}
$y = (x - 1)^4$.
\end{prob}

\begin{prob}[S]{1417}
$y = x^m (1 - x)^n$（$m$ 及 $n$ 为正整数）.
\end{prob}

\begin{prob}[S]{1418}
$y = \cos x + \cosh x$.
\end{prob}

\begin{prob}[S]{1419}
$y = (x + 1)^{10} \mathrm{e}^{-x}$.
\end{prob}

\begin{prob}[S]{1420}
$y = \left(1 + x + \dfrac{x^2}{2!} + \cdots + \dfrac{x^n}{n!}\right)\mathrm{e}^{-x}$（$n$ 为正整数）
\end{prob}

\begin{prob}[S]{1421}
$y = |x|$.
\end{prob}

\begin{prob}[S]{1422}
$y = x^{\frac{1}{3}} (1 - x)^{\frac{2}{3}}$.
\end{prob}

\begin{prob}[M]{1423}
设
\[
f(x) = (x - x_0)^n \varphi(x) \quad (n \text{ 为正整数}),
\]
其中函数 $\varphi(x)$ 当 $x = x_0$ 时连续，且 $\varphi(x_0) \neq 0$. 研究此函数在点 $x = x_0$ 的极值.
\end{prob}

\begin{prob}[M]{1424}
设
\[
f(x) = \frac{P(x)}{Q(x)}, \quad f'(x) = \frac{P_1(x)}{Q^2(x)},
\]
$x_0$ 为函数 $f(x)$ 的临界点，即
\[
P_1(x_0) = 0, \quad Q(x_0) \neq 0.
\]
证明：
\[
\operatorname{sgn} f''(x_0) = \operatorname{sgn} P_1'(x_0).
\]
\end{prob}

\begin{prob}[M]{1425}
可否断定：若函数 $f(x)$ 在点 $x_0$ 有极大值，则在此点的某充分小邻域内，函数 $f(x)$ 在点 $x_0$ 的左侧递增，而在其右侧递减? 研究一个例子：
\[
f(x) = \begin{cases}
2 - x^2 \left(2 + \sin\dfrac{1}{x}\right), & x \neq 0, \\
2, & x = 0.
\end{cases}
\]
\end{prob}

\begin{prob}[L]{1426}
证明：函数
\[
f(x) = \begin{cases}
\mathrm{e}^{-\frac{1}{x^2}}, & x \neq 0, \\
0, & x = 0
\end{cases}
\]
在点 $x = 0$ 有极小值，而函数
\[
g(x) = \begin{cases}
x\mathrm{e}^{-\frac{1}{x^2}}, & x \neq 0, \\
0, & x = 0
\end{cases}
\]
在点 $x = 0$ 没有极值，尽管
\[
f^{(n)}(0) = 0, \quad g^{(n)}(0) = 0 \quad (n = 1, 2, \cdots).
\]
作出这些函数的图像.
\end{prob}

\begin{prob}[L]{1427}
研究下列函数的极值并作出其图像：
\begin{enumerate}[label=(\alph*), itemsep=2pt, topsep=3pt]
    \item $f(x) = \begin{cases} \mathrm{e}^{-\frac{1}{|x|}}\left(\sqrt{2} + \sin\dfrac{1}{x}\right), & x \neq 0, \\ 0, & x = 0\text{；} \end{cases}$
    \item $f(x) = \begin{cases} \mathrm{e}^{-\frac{1}{|x|}}\left(\sqrt{2} + \cos\dfrac{1}{x}\right), & x \neq 0, \\ 0, & x = 0. \end{cases}$
\end{enumerate}
\end{prob}

\begin{prob}[M]{1428}
研究函数
\[
f(x) = \begin{cases}
|x|\left(2 + \cos\dfrac{1}{x}\right), & x \neq 0, \\
0, & x = 0
\end{cases}
\]
在点 $x = 0$ 处的极值，并作出此函数的图像.
\end{prob}

\probgroup{求下列函数的极值：}

\begin{prob}[S]{1429}
$y = x^3 - 6x^2 + 9x - 4$.
\end{prob}

\begin{prob}[S]{1430}
$y = 2x^2 - x^4$.
\end{prob}

\begin{prob}[S]{1431}
$y = x(x - 1)^2 (x - 2)^3$.
\end{prob}

\begin{prob}[S]{1432}
$y = x + \dfrac{1}{x}$.
\end{prob}

\begin{prob}[S]{1433}
$y = \dfrac{2x}{1 + x^2}$.
\end{prob}

\begin{prob}[S]{1434}
$y = \dfrac{x^2 - 3x + 2}{x^2 + 2x + 1}$.
\end{prob}

\begin{prob}[S]{1435}
$y = \sqrt{2x - x^2}$.
\end{prob}

\begin{prob}[S]{1436}
$y = x\sqrt[3]{x - 1}$.
\end{prob}

\begin{prob}[S]{1437}
$y = x\mathrm{e}^{-x}$.
\end{prob}

\begin{prob}[S]{1438}
$y = \sqrt{x} \ln x$.
\end{prob}

\begin{prob}[S]{1439}
$y = \dfrac{\ln^2 x}{x}$.
\end{prob}

\begin{prob}[S]{1440}
$y = \cos x + \dfrac{1}{2} \cos 2x$.
\end{prob}

\begin{prob}[S]{1441}
$y = \dfrac{10}{1 + \sin^2 x}$.
\end{prob}

\begin{prob}[S]{1442}
$y = \arctan x - \dfrac{1}{2} \ln(1 + x^2)$.
\end{prob}

\begin{prob}[S]{1443}
$y = \mathrm{e}^x \sin x$.
\end{prob}

\begin{prob}[S]{1444}
$y = |x|\mathrm{e}^{-|x - 1|}$.
\end{prob}

\probgroup{求下列函数在所给闭区间上的最大值和最小值：}

\begin{prob}[S]{1445}
$f(x) = 2^x, \ [-1, 5]$.
\end{prob}

\begin{prob}[S]{1446}
$f(x) = x^2 - 4x + 6, \ [-3, 10]$.
\end{prob}

\begin{prob}[S]{1447}
$f(x) = |x^2 - 3x + 2|, \ [-10, 10]$.
\end{prob}

\begin{prob}[S]{1448}
$f(x) = x + \dfrac{1}{x}, \ [0.01, 100]$.
\end{prob}

\begin{prob}[S]{1449}
$f(x) = \sqrt{5 - 4x}, \ [-1, 1]$.
\end{prob}

\probgroup{求下列函数在所给区间上的下确界 ($\inf$) 与上确界 ($\sup$)：}

\begin{prob}[S]{1450}
$f(x) = x\mathrm{e}^{-0.01x}, \ (0, +\infty)$.
\end{prob}

\begin{prob}[S]{1451}
$f(x) = \left(1 + x + \dfrac{x^2}{2!} + \cdots + \dfrac{x^n}{n!}\right)\mathrm{e}^{-x}, \ (0, +\infty)$.
\end{prob}

\begin{prob}[S]{1452}
$f(x) = \dfrac{1 + x^2}{1 + x^4}, \ (0, +\infty)$.
\end{prob}

\begin{prob}[S]{1453}
$f(x) = \mathrm{e}^{-x^2} \cos x^2, \ (-\infty, +\infty)$.
\end{prob}

\begin{prob}[M]{1454.1}
求函数 $f(\xi) = \dfrac{1 + \xi}{3 + \xi^2}$ 在区间 $x < \xi < +\infty$ 内的下确界与上确界. 作出下列函数的图像：
\[
M(x) = \sup_{x < \xi < +\infty} f(\xi), \quad m(x) = \inf_{x < \xi < +\infty} f(\xi).
\]
\end{prob}

\begin{prob}[M]{1454.2}
设
\[
M_k = \sup_x |f^{(k)}(x)| \quad (k = 0, 1, 2, \cdots).
\]
若 $f(x) = \mathrm{e}^{-x^2}$，求 $M_0, M_1$ 和 $M_2$.
\end{prob}

\begin{prob}[L]{1455}
求以下各数列的最大项：
\begin{enumerate}[label=(\alph*), itemsep=2pt, topsep=3pt]
    \item $\dfrac{n^{10}}{2^n} \ (n = 1, 2, \cdots)$；
    \item $\dfrac{\sqrt{n}}{n + 10\,000} \ (n = 1, 2, \cdots)$；
    \item $\sqrt[n]{n} \ (n = 1, 2, \cdots)$.
\end{enumerate}
\end{prob}

\begin{prob}[L]{1456.1}
证明下列不等式：
\begin{enumerate}[label=(\alph*), itemsep=2pt, topsep=3pt]
    \item $|3x - x^3| \leqslant 2 \ (|x| \leqslant 2)$；
    \item $\dfrac{1}{2^{p - 1}} \leqslant x^p + (1 - x)^p \leqslant 1 \ (0 \leqslant x \leqslant 1, \, p > 1)$；
    \item $x^m (a - x)^n \leqslant \dfrac{m^m n^n}{(m + n)^{m + n}} a^{m + n} \ (0 \leqslant x \leqslant a, \, m > 0, \, n > 0)$；
    \item $\dfrac{x + a}{2^{\frac{n - 1}{n}}} \leqslant \sqrt[n]{x^n + a^n} \leqslant x + a \ (x > 0, \, a > 0, \, n > 1)$；
    \item $|a\sin x + b\cos x| \leqslant \sqrt{a^2 + b^2}$.
\end{enumerate}
\end{prob}

\begin{prob}[M]{1456.2}
证明不等式：
\[
\frac{2}{3} \leqslant \frac{x^2 + 1}{x^2 + x + 1} \leqslant 2 \quad (-\infty < x < +\infty).
\]
\end{prob}

\begin{prob}[M]{1457}
求多项式
\[
P(x) = x(x - 1)^2 (x + 2)
\]
在闭区间 $[-2, 1]$ 上``与零的偏差''，就是求
\[
E_P = \sup_{-2 \leqslant x \leqslant 1} |P(x)|.
\]
\end{prob}

\begin{prob}[M]{1458}
应当选择怎样的系数 $q$，使多项式
\[
P(x) = x^2 + q
\]
在闭区间 $[-1, 1]$ 上与零的偏差最小，即
\[
E_P = \sup_{-1 \leqslant x \leqslant 1} |P(x)| = \min ?
\]
\end{prob}

\begin{prob}[M]{1459}
数
\[
\Delta = \sup_{a \leqslant x \leqslant b} |f(x) - g(x)|
\]
称为函数 $f(x)$ 与 $g(x)$ 在闭区间 $[a, b]$ 上的绝对偏差.

求函数 $f(x) = x^2$ 和 $g(x) = x^3$ 在闭区间 $[0, 1]$ 上的绝对偏差.
\end{prob}

\begin{prob}[M]{1460}
在闭区间 $[x_1, x_2]$ 上用线性函数
\[
g(x) = (x_1 + x_2)x + b
\]
近似地代替函数
\[
f(x) = x^2,
\]
使函数 $f(x)$ 与 $g(x)$ 的绝对偏差（参阅上题）最小，并求此最小的绝对偏差.
\end{prob}

\begin{prob}[M]{1461}
求函数
\[
f(x) = \max\{2|x|, |1+x|\}
\]
的极小值.
\end{prob}

\probgroup{确定下列各方程实根的数目，并划分出这些根所在的区间：}

\begin{prob}[S]{1462}
$x^3 - 6x^2 + 9x - 10 = 0.$
\end{prob}

\begin{prob}[S]{1463}
$x^3 - 3x^2 - 9x + h = 0.$
\end{prob}

\begin{prob}[S]{1464}
$3x^4 - 4x^3 - 6x^2 + 12x - 20 = 0.$
\end{prob}

\begin{prob}[S]{1465}
$x^5 - 5x = a.$
\end{prob}

\begin{prob}[S]{1466}
$\ln x = kx.$
\end{prob}

\begin{prob}[S]{1467}
$\mathrm{e}^x = ax^2.$
\end{prob}

\begin{prob}[S]{1468}
当 $0 \leqslant x \leqslant \pi$ 时，$\sin^3 x \cdot \cos x = a.$
\end{prob}

\begin{prob}[S]{1469}
$\cosh x = kx.$
\end{prob}

\begin{prob}[M]{1470}
在什么条件下方程
\[
x^3 + px + q = 0
\]
有：(a) 一个实根；(b) 三个实根？在平面 $(p, q)$ 上画出相应的区域.
\end{prob}

\probsec{§12. 依据函数的特征点作函数图像}

\probgroup{作出下列函数的图像：}

\begin{prob}[S]{1471}
$y = 3x - x^3.$
\end{prob}

\begin{prob}[S]{1472}
$y = 1 + x^2 - \frac{x^4}{2}.$
\end{prob}

\begin{prob}[S]{1473}
$y = (x+1)(x-2)^2.$
\end{prob}

\begin{prob}[S]{1474*}
$y = \frac{2 - x^2}{1 + x^4}.$
\end{prob}

\begin{prob}[S]{1475*}
$y = \frac{x^2 - 1}{x^2 - 5x + 6}.$
\end{prob}

\begin{prob}[S]{1476*}
$y = \frac{x}{(1+x)(1-x)^2}.$
\end{prob}

\begin{prob}[S]{1477}
$y = \frac{x^4}{(1+x)^3}.$
\end{prob}

\begin{prob}[S]{1478}
$y = \left(\frac{1+x}{1-x}\right)^4.$
\end{prob}

\begin{prob}[S]{1479}
$y = \frac{x^2(x-1)}{(x+1)^2}.$
\end{prob}

\begin{prob}[S]{1480}
$y = \frac{x}{(1-x^2)^2}.$
\end{prob}

\begin{prob}[S]{1481}
$y = \frac{(x+1)^3}{(x-1)^2}.$
\end{prob}

\begin{prob}[S]{1482*}
$y = \frac{x^4 + 8}{x^3 + 1}.$
\end{prob}

\begin{prob}[S]{1483}
$y = \frac{1}{1+x} - \frac{10}{3x^2} + \frac{1}{1-x}.$
\end{prob}

\begin{prob}[S]{1484}
$y = (x-3)\sqrt{x}.$
\end{prob}

\begin{prob}[M]{1485}
(a) $y = \pm\sqrt{8x^2 - x^4}$; \qquad (b) $y = \frac{x-2}{\sqrt{x^2+1}}.$
\end{prob}

\begin{prob}[S]{1486}
$y = \pm\sqrt{(x-1)(x-2)(x-3)}.$
\end{prob}

\begin{prob}[S]{1487*}
$y = \sqrt[3]{x^3 - x^2 - x + 1}.$
\end{prob}

\begin{prob}[S]{1488}
$y = \sqrt[3]{x^2} - \sqrt[3]{x^2 + 1}.$
\end{prob}

\begin{prob}[S]{1489}
$y = (x+2)^{\frac{2}{3}} - (x-2)^{\frac{2}{3}}.$
\end{prob}

\begin{prob}[S]{1490}
$y = (x+1)^{\frac{2}{3}} + (x-1)^{\frac{2}{3}}.$
\end{prob}

\begin{prob}[S]{1491}
$y = \frac{x}{\sqrt[3]{x^2 - 1}}.$
\end{prob}

\begin{prob}[S]{1492}
$y = \frac{x^2\sqrt{x^2 - 1}}{2x^2 - 1}.$
\end{prob}

\begin{prob}[S]{1493}
$y = \frac{|1+x|^{\frac{3}{2}}}{\sqrt{x}}.$
\end{prob}

\begin{prob}[S]{1494}
$y = 1 - x + \sqrt{\frac{x^3}{3+x}}.$
\end{prob}

\begin{prob}[S]{1495}
$y = \sqrt[3]{\frac{x^2}{x+1}}.$
\end{prob}

\begin{prob}[S]{1496*}
$y = \sqrt{\frac{x^4 + 3}{x^2 + 1}}.$
\end{prob}

\begin{prob}[S]{1497}
$y = \sin x + \cos^2 x.$
\end{prob}

\begin{prob}[S]{1498}
$y = (7 + 2\cos x)\sin x.$
\end{prob}

\begin{prob}[S]{1499}
$y = \sin x + \frac{1}{3}\sin 3x.$
\end{prob}

\begin{prob}[S]{1500}
$y = \cos x - \frac{1}{2}\cos 2x.$
\end{prob}

\begin{prob}[S]{1501}
$y = \sin^4 x + \cos^4 x.$
\end{prob}

\begin{prob}[S]{1502}
$y = \sin x \cdot \sin 3x.$
\end{prob}

\begin{prob}[S]{1503}
$y = \frac{\sin x}{\sin\left(x + \dfrac{\pi}{4}\right)}.$
\end{prob}

\begin{prob}[M]{1504}
(a) $y = \frac{\cos x}{\cos 2x}$; \qquad (b) $y = \frac{\sin x}{2 + \cos x}.$
\end{prob}

\begin{prob}[S]{1505}
$y = 2x - \tan x.$
\end{prob}

\begin{prob}[S]{1506}
$y = \mathrm{e}^{2x - x^2}.$
\end{prob}

\begin{prob}[S]{1507}
$y = (1 + x^2)\mathrm{e}^{-x^2}.$
\end{prob}

\begin{prob}[S]{1508}
$y = x + \mathrm{e}^{-x}.$
\end{prob}

\begin{prob}[M]{1509}
(a) $y = x^{\frac{2}{3}}\mathrm{e}^{-x}$; \qquad (b) $y = \mathrm{e}^{-2x}\sin^2 x.$
\end{prob}

\begin{prob}[S]{1510}
$y = \frac{\mathrm{e}^x}{1+x}.$
\end{prob}

\begin{prob}[S]{1511}
$y = \sqrt{1 - \mathrm{e}^{-x^2}}.$
\end{prob}

\begin{prob}[S]{1512}
$y = \frac{\ln x}{\sqrt{x}}.$
\end{prob}

\begin{prob}[S]{1513}
$y = \ln(x + \sqrt{x^2 + 1}).$
\end{prob}

\begin{prob}[S]{1514}
$y = \sqrt{x^2 + 1}\cdot\ln(x + \sqrt{x^2 + 1}).$
\end{prob}

\begin{prob}[S]{1515}
$y = \frac{\arcsin x}{\sqrt{1 - x^2}}.$
\end{prob}

\begin{prob}[S]{1516}
$y = x + \arctan x.$
\end{prob}

\begin{prob}[S]{1517}
$y = \frac{x}{2} + \operatorname{arccot} x.$
\end{prob}

\begin{prob}[S]{1518}
$y = x\arctan x.$
\end{prob}

\begin{prob}[S]{1519}
$y = \arcsin\frac{2x}{1 + x^2}.$
\end{prob}

\begin{prob}[S]{1520}
$y = \arccos\frac{1 - x^2}{1 + x^2}.$
\end{prob}

\begin{prob}[S]{1521}
$y = (x+2)\mathrm{e}^{\frac{1}{x}}.$
\end{prob}

\begin{prob}[S]{1522}
$y = 2^{\sqrt{x^2+1} - \sqrt{x^2-1}}.$
\end{prob}

\begin{prob}[S]{1523*}
$y = \ln\frac{x^2 - 3x + 2}{x^2 + 1}.$
\end{prob}

\begin{prob}[S]{1524}
$y = a\arcsin\frac{x}{a} - \sqrt{a^2 - x^2}\ (a > 0).$
\end{prob}

\begin{prob}[S]{1525}
$y = \arccos\frac{1 - x}{1 - 2x}.$
\end{prob}

\begin{prob}[S]{1526}
$y = x^x.$
\end{prob}

\begin{prob}[S]{1527}
$y = x^{\frac{1}{x}}.$
\end{prob}

\begin{prob}[S]{1528}
$y = (1 + x)^{\frac{1}{x}}.$
\end{prob}

\begin{prob}[S]{1529*}
$y = x\left(1 + \frac{1}{x}\right)^x\ (x > 0).$
\end{prob}

\begin{prob}[S]{1530*}
$y = \frac{\mathrm{e}^{\frac{1}{1 - x^2}}}{1 + x^2}$ (不研究凹凸性).
\end{prob}

\probgroup{作出下列参数方程所表示的曲线：}

\begin{prob}[S]{1531}
$x = \frac{(t+1)^2}{4},\ y = \frac{(t-1)^2}{4}.$
\end{prob}

\begin{prob}[S]{1532}
$x = 2t - t^2,\ y = 3t - t^3.$
\end{prob}

\begin{prob}[S]{1533*}
$x = \frac{t^2}{t - 1},\ y = \frac{t}{t^2 - 1}.$
\end{prob}

\begin{prob}[S]{1534}
$x = \frac{t^2}{1 - t^2},\ y = \frac{1}{1 + t^2}.$
\end{prob}

\begin{prob}[S]{1535}
$x = t + \mathrm{e}^{-t},\ y = 2t + \mathrm{e}^{-2t}.$
\end{prob}

\begin{prob}[S]{1536}
$x = a\cos 2t,\ y = a\cos 3t\ (a > 0).$
\end{prob}

\begin{prob}[S]{1537}
$x = \cos^4 t,\ y = \sin^4 t.$
\end{prob}

\begin{prob}[S]{1538}
$x = t\ln t,\ y = \frac{\ln t}{t}.$
\end{prob}

\begin{prob}[S]{1539}
$x = \frac{a}{\cos^3 t},\ y = a\tan^3 t\ (a > 0).$
\end{prob}

\begin{prob}[S]{1540}
$x = a(\sinh t - t),\ y = a(\cosh t - 1)\ (a > 0).$
\end{prob}

\probgroup{把下列曲线方程化为参数方程，然后作出这些曲线的图像：}

\begin{prob}[M]{1541}
$x^3 + y^3 - 3axy = 0\ (a > 0).$
\probhint{令 $y = tx$.}
\end{prob}

\begin{prob}[S]{1542}
$x^2 + y^2 = x^4 + y^4.$
\end{prob}

\begin{prob}[S]{1543}
$x^2 y^2 = x^3 - y^3.$
\end{prob}

\begin{prob}[S]{1544}
$x^y = y^x\ (x > 0, y > 0).$
\end{prob}

\begin{prob}[S]{1545}
作出曲线 $\cosh^2 x - \cosh^2 y = 1$ 的图像.
\end{prob}

\probgroup{作出下列用极坐标 $(\varphi, r)\ (r \geqslant 0)$ 表示的函数的图像：}

\begin{prob}[S]{1546}
$r = a + b\cos\varphi\ (0 < a \leqslant b).$
\end{prob}

\begin{prob}[S]{1547}
$r = a\sin 3\varphi\ (a > 0).$
\end{prob}

\begin{prob}[S]{1548}
$r = \frac{a}{\sqrt{\cos 3\varphi}}\ (a > 0).$
\end{prob}

\begin{prob}[S]{1549*}
$r = a\frac{\tanh\varphi}{\varphi - 1}$，其中 $\varphi > 1\ (a > 0).$
\end{prob}

\begin{prob}[S]{1550*}
$\varphi = \arccos\frac{r - 1}{r^2}.$
\end{prob}

\probgroup{作出下列曲线族的图像（$a$ 为参变量）：}

\begin{prob}[S]{1551}
$y = x^2 - 2x + a.$
\end{prob}

\begin{prob}[S]{1552}
$y = x + \frac{a^2}{x}.$
\end{prob}

\begin{prob}[S]{1553}
$y = x \pm \sqrt{a(1 - x^2)}.$
\end{prob}

\begin{prob}[S]{1554}
$y = \frac{x}{2} + \mathrm{e}^{-ax}.$
\end{prob}

\begin{prob}[S]{1555}
$y = x\mathrm{e}^{-\frac{x}{a}}.$
\end{prob}

\probsec{§13. 函数的极大值与极小值问题}

\begin{prob}[M]{1556}
证明：若函数 $f(x)$ 不为负，则函数
\[
F(x) = Cf^2(x)\ (C > 0)
\]
与函数 $f(x)$ 有相同的极值点.
\end{prob}

\begin{prob}[M]{1557}
证明：若当 $-\infty < x < +\infty$ 时函数 $\varphi(x)$ 严格单调递增，则函数
\[
f(x) \quad \text{和} \quad \varphi(f(x))
\]
有相同的极值点.
\end{prob}

\begin{prob}[M]{1558}
若二正数之和等于常数 $a$，求此二数的 $m$ 次幂与 $n$ 次幂 $(m > 0, n > 0)$ 之积的最大值.
\end{prob}

\begin{prob}[M]{1559}
若二正数之积等于常数 $a$，求此二数的 $m$ 次幂与 $n$ 次幂 $(m > 0, n > 0)$ 之和的最小值.
\end{prob}

\begin{prob}[M]{1560}
当对数之底取何值时存在这样的数，它本身和它的对数相等？
\end{prob}

\begin{prob}[M]{1561}
从面积为给定值 $S$ 的一切矩形中，求其周长为最小者.
\end{prob}

\begin{prob}[M]{1562}
若直角三角形的一直角边与斜边之和为常数，求具有最大面积的直角三角形.
\end{prob}

\begin{prob}[M]{1563}
要使容积为给定值 $V$ 的圆柱形闭合容器具有最小的表面积，其尺寸如何？
\end{prob}

\begin{prob}[M]{1564}
在不超过半圆的给定弓形内作出具有最大面积的内接矩形.
\end{prob}

\begin{prob}[M]{1565}
在椭圆
\[
\frac{x^2}{a^2} + \frac{y^2}{b^2} = 1
\]
内作出具有最大面积而边平行于椭圆轴的内接矩形.
\end{prob}

\begin{prob}[M]{1566}
在底边为 $b$ 及高为 $h$ 的三角形内作出具有最大周长的内接矩形. 研究此问题有解的可能性.
\end{prob}

\begin{prob}[M]{1567}
从直径为 $d$ 的原木（圆形横截面）切出横截面为矩形的梁，此矩形的底等于 $b$，高等于 $h$. 若梁的强度\footnote{此处的强度指梁的截面模量. —— 译注}与 $bh^2$ 成正比，问梁的尺寸如何其强度最大？
\end{prob}

\begin{prob}[M]{1568}
在半径为 $R$ 的半球中作出具有最大体积且底为正方形的内接长方体.
\end{prob}

\begin{prob}[M]{1569}
在半径为 $R$ 的球内作出具有最大体积的内接圆柱体.
\end{prob}

\begin{prob}[M]{1570}
在半径为 $R$ 的球内作出具有最大表面积的内接圆柱体.
\end{prob}

\begin{prob}[M]{1571}
在已知球外作出具有最小体积的外切圆锥体.
\end{prob}

\begin{prob}[M]{1572}
求母线为给定值 $l$ 的圆锥之最大体积.
\end{prob}

\begin{prob}[M]{1573}
在顶角为 $2\alpha$ 底半径为 $R$ 的直圆锥体中作出具有最大表积的内接圆柱体.
\end{prob}

\begin{prob}[M]{1574}
求从点 $M(p, p)$ 到抛物线 $y^2 = 2px$ 的最短距离.
\end{prob}

\begin{prob}[M]{1575}
求从点 $A(2, 0)$ 到圆 $x^2 + y^2 = 1$ 的最短与最长距离.
\end{prob}

\begin{prob}[M]{1576}
求椭圆 $\dfrac{x^2}{a^2} + \dfrac{y^2}{b^2} = 1\ (0 < b < a)$ 的经过顶点 $B(0, -b)$ 的最大弦.
\end{prob}

\begin{prob}[M]{1577}
过椭圆 $\dfrac{x^2}{a^2} + \dfrac{y^2}{b^2} = 1$ 上的点 $M(x, y)$ 引切线，使由此切线与坐标轴构成的三角形具有最小的面积.
\end{prob}

\begin{prob}[M]{1578}
直径相同的圆柱体与半球体拼接在一起构成一物体，其体积为 $V$. 要使此物体具有最小的表面积，其尺寸如何？
\end{prob}

\begin{prob}[M]{1579}
设明渠的横截面为等腰梯形，渠中流水的横截面面积为 $S$，水面的高为 $h$. 问水渠侧边的倾角 $\varphi$ 如何，才使其横截面边界被水浸湿的部分具有最小的长度？
\end{prob}

\begin{prob}[M]{1580}
设封闭曲线所围面积为 $S$，则该曲线的周长与面积同为 $S$ 的圆的周长之比称为封闭曲线的弯曲度.

设等腰梯形 $ABCD\ (AD // BC)$ 的底边 $AD = 2a$，锐角 $BAD = \alpha$，问等腰梯形的形状如何，才有最小的弯曲度？
\end{prob}

\begin{prob}[M]{1581}
从半径为 $R$ 的圆中应切去怎样的扇形，才能使余下的部分可卷成一漏斗，其容积为最大？
\end{prob}

\begin{prob}[M]{1582}
从南至北的铁路经过 $B$ 城，工厂 $A$ 距此铁路的最短距离为 $a$，距偏北方的 $B$ 城所在纬线的距离为 $b$. 为了从 $A$ 至 $B$ 运输货物最经济，从工厂修建一条专线铁路. 若每吨货物沿专线运输的价格是 $p\ \text{元/km}$，而沿常规铁路为 $q\ \text{元/km}\ (p > q)$，则支线应向常规铁路取怎样的角度 $\varphi$？
\end{prob}

\begin{prob}[M]{1583}
二船各以恒定的速度 $u$ 和 $v$ 沿直线前进，二者前进方向所成的角为 $\theta$. 若在某时刻它们到其路线交点的距离分别为 $a$ 和 $b$，求二船的最小距离.
\end{prob}

\begin{prob}[M]{1584}
在 $A$ 与 $B$ 二点处各有一光源，其发光强度分别为 $S_1$ 与 $S_2$. 在线段 $AB = a$ 上求出最小照度的点 $M$.
\end{prob}

\begin{prob}[M]{1585}
发光点位于半径为 $R$ 与 $r\ (R > r)$ 的二互不相交之球的连心线上，并在此二球的外面. 此发光点的位置如何，才可使二球表面上被照明部分之和为最大？
\end{prob}

\begin{prob}[M]{1586}
设圆桌面的半径为 $a$. 应当在正对桌面中央多高的地方安置电灯，才可使桌面边缘的照度为最大？
\probhint{照度可用下列公式表示：
\[
I = I_0\frac{\sin\varphi}{r^2},
\]
式中 $\varphi$ 为光线对桌面的倾角，$r$ 为光源与被照面的距离，$I_0$ 为光源的发光强度.}
\end{prob}

\begin{prob}[M]{1587}
向宽为 $a$ 的河修建一宽为 $b$ 的运河，二者成直角相交，问能驶进这运河的船，其最大的长度如何？
\end{prob}

\begin{prob}[M]{1588}
船航行一昼夜的耗费由两部分组成：固定部分等于 $a$ 元，变动部分与速度的立方成正比增加. 在怎样的速度 $v$ 时，航行最为经济？
\end{prob}

\begin{prob}[M]{1589}
重量为 $P$ 的物体位于粗糙的水平面上，需用力把物体从原位置移动. 若物体摩擦因子等于 $k$，问作用力对水平面的倾斜程度如何，才使所需的力为最小？
\end{prob}

\begin{prob}[M]{1590}
有一茶杯，其形状是半径为 $a$ 的半球，在茶杯中放一长为 $l > 2a$ 的杆，求杆的平衡位置.
\end{prob}

\probsec{§14. 曲线的相切. 曲率圆. 渐屈线}

\begin{prob}[M]{1591}
选择直线
\[
y = kx + b
\]
的参数 $k$ 与 $b$，使它与曲线
\[
y = x^3 - 3x^2 + 2
\]
二阶或更高阶相切.
\end{prob}

\begin{prob}[M]{1592}
应该怎样选择系数 $a, b$ 和 $c$，才能使抛物线
\[
y = ax^2 + bx + c
\]
在点 $x = x_0$ 与曲线 $y = \mathrm{e}^x$ 二阶相切？
\end{prob}

\begin{prob}[L]{1593}
下列曲线与 $Ox$ 轴在点 $x = 0$ 相切的阶如何：
\begin{enumerate}[label=(\alph*), itemsep=2pt, topsep=3pt]
\item $y = 1 - \cos x$;
\item $y = \tan x - \sin x$;
\item $y = \mathrm{e}^x - \left(1 + x + \frac{x^2}{2}\right)$.
\end{enumerate}
\end{prob}

\begin{prob}[M]{1594}
证明：曲线
\[
y = \begin{cases} \mathrm{e}^{-\frac{1}{x^2}}, & x \neq 0, \\ 0, & x = 0 \end{cases}
\]
在点 $x = 0$ 与 $Ox$ 轴相切的阶为无穷大.
\end{prob}

\begin{prob}[M]{1595}
求双曲线
\[
xy = 1
\]
在下列各点的曲率半径和曲率中心：
(a) $M(1, 1)$; \qquad (b) $N(100, 0.01)$.
\end{prob}

\probgroup{求下列曲线的曲率半径：}

\begin{prob}[S]{1596}
抛物线 $y^2 = 2px$.
\end{prob}

\begin{prob}[S]{1597}
椭圆 $\frac{x^2}{a^2} + \frac{y^2}{b^2} = 1 \ (a \geqslant b > 0)$.
\end{prob}

\begin{prob}[S]{1598}
双曲线 $\frac{x^2}{a^2} - \frac{y^2}{b^2} = 1$.
\end{prob}

\begin{prob}[S]{1599}
星形线 $x^{\frac{2}{3}} + y^{\frac{2}{3}} = a^{\frac{2}{3}}$.
\end{prob}

\begin{prob}[S]{1600}
椭圆 $x = a\cos t,\ y = b\sin t$.
\end{prob}

\begin{prob}[S]{1601}
摆线 $x = a(t - \sin t),\ y = a(1 - \cos t)$.
\end{prob}

\begin{prob}[S]{1602}
圆的渐伸线 $x = a(\cos t + t\sin t),\ y = a(\sin t - t\cos t)$.
\end{prob}

\begin{prob}[M]{1603}
证明：二次曲线
\[
y^2 = 2px - qx^2
\]
的曲率半径与法线段的立方成正比.
\end{prob}

\begin{prob}[S]{1604}
写出以极坐标表示的曲线的曲率半径公式.
\end{prob}

\probgroup{求下列极坐标方程所表示的曲线的曲率半径：}

\begin{prob}[S]{1605}
阿基米德螺线 $r = a\varphi$.
\end{prob}

\begin{prob}[S]{1606}
对数螺线 $r = a\mathrm{e}^{m\varphi}$.
\end{prob}

\begin{prob}[S]{1607}
心脏线 $r = a(1 + \cos\varphi)$.
\end{prob}

\begin{prob}[S]{1608}
双纽线 $r^2 = a^2 \cos 2\varphi$.
\end{prob}

\begin{prob}[M]{1609}
在曲线 $y = \ln x$ 上求曲率最大的点.
\end{prob}

\begin{prob}[M]{1610}
三次抛物线 $y = \frac{kx^3}{6} \ (0 \leqslant x < +\infty, k > 0)$ 的最大曲率等于 $\frac{1}{1000}$，求达到此最大曲率的点 $x$.
\end{prob}

\probgroup{求下列各曲线的渐屈线方程：}

\begin{prob}[S]{1611}
抛物线 $y^2 = 2px$.
\end{prob}

\begin{prob}[S]{1612}
椭圆 $\frac{x^2}{a^2} + \frac{y^2}{b^2} = 1$.
\end{prob}

\begin{prob}[S]{1613}
星形线 $x^{\frac{2}{3}} + y^{\frac{2}{3}} = a^{\frac{2}{3}}$.
\end{prob}

\begin{prob}[S]{1614}
曳物线 $x = a\ln\frac{a + \sqrt{a^2 - y^2}}{y} - \sqrt{a^2 - y^2}$.
\end{prob}

\begin{prob}[S]{1615}
对数螺线 $r = a\mathrm{e}^{m\varphi}$.
\end{prob}

\begin{prob}[M]{1616}
证明：摆线
\[
x = a(t - \sin t),\quad y = a(1 - \cos t)
\]
的渐屈线仍为一摆线，仅其位置与所给摆线不同而已.
\end{prob}

\probsec{§15. 方程的近似解法}

\probgroup{利用比例法求下列方程的根（精确到 $0.001$）：}

\begin{prob}[S]{1617}
$x^3 - 6x + 2 = 0$.
\end{prob}

\begin{prob}[S]{1618}
$x^4 - x - 1 = 0$.
\end{prob}

\begin{prob}[S]{1619}
$x - 0.1\sin x = 2$.
\end{prob}

\begin{prob}[S]{1620}
$\cos x = x^2$.
\end{prob}

\probgroup{利用牛顿法求下列方程的根（精确到所指定的精度）：}

\begin{prob}[S]{1621}
$x^2 + \frac{1}{x^2} = 10x$（精确到 $10^{-3}$）.
\end{prob}

\begin{prob}[S]{1622}
$x\lg x = 1$（精确到 $10^{-4}$）.
\end{prob}

\begin{prob}[S]{1623}
$\cos x \cdot \cosh x = 1$（精确到 $10^{-3}$）（二正根）.
\end{prob}

\begin{prob}[S]{1624}
$x + \mathrm{e}^x = 0$（精确到 $10^{-5}$）.
\end{prob}

\begin{prob}[S]{1625}
$x\tanh x = 1$（精确到 $10^{-6}$）.
\end{prob}

\begin{prob}[M]{1626}
求方程
\[
\tan x = x
\]
最小的三个正根（精确到 $0.001$）.
\end{prob}

\begin{prob}[M]{1627}
求方程
\[
\cot x = \frac{1}{x} - \frac{x}{2}
\]
的二正根（精确到 $10^{-3}$）.
\end{prob}
